%% Ngân hàng đề Toán 12 (KNTT) -- Chương IV, Bài 12, Dạng 01
%% Dạng: Tính chất của tích phân, tích phân theo định nghĩa
%% Phương án ĐẦU TIÊN trong \cauTN luôn là đáp án đúng (gói trondethi tự đảo).
%% Phần II: \y{nội dung}{d|s}. Phần III: đáp án đặt ở đối số cuối.

\cauTN[id=C4-B12-D01-NB-001,muc=NB,loigiai={$\displaystyle\int_{1}^{3}f(x)\,\mathrm{d}x=F(3)-F(1)=4-(-4)=8$.}]{Cho hàm số $f(x)$ liên tục trên $\mathbb{R}$ và $F(x)$ là một nguyên hàm của $f(x)$, thoả mãn $F(1)=-4$, $F(3)=4$. Giá trị của $\displaystyle\int_{1}^{3}f(x)\,\mathrm{d}x$ bằng}
  {$8$}
  {$-8$}{$0$}{$-16$}

\cauTN[id=C4-B12-D01-NB-002,muc=NB,loigiai={$\displaystyle\int_{-1}^{0}\left[-4f(x)-3g(x)\right]\,\mathrm{d}x=-4\cdot \displaystyle\int_{-1}^{0}f(x)\,\mathrm{d}x-3\cdot \displaystyle\int_{-1}^{0}g(x)\,\mathrm{d}x=-1$.}]{Biết $\displaystyle\int_{-1}^{0}f(x)\,\mathrm{d}x=1$ và $\displaystyle\int_{-1}^{0}g(x)\,\mathrm{d}x=-1$. Giá trị của $\displaystyle\int_{-1}^{0}\left[-4f(x)-3g(x)\right]\,\mathrm{d}x$ bằng}
  {$-1$}
  {$-7$}{$0$}{$1$}

\cauTN[id=C4-B12-D01-NB-003,muc=NB,loigiai={$\displaystyle\int_{-1}^{-2}f(x)\,\mathrm{d}x=-\displaystyle\int_{-2}^{-1}f(x)\,\mathrm{d}x=-5$.}]{Biết $\displaystyle\int_{-2}^{-1}f(x)\,\mathrm{d}x=5$. Giá trị của $\displaystyle\int_{-1}^{-2}f(x)\,\mathrm{d}x$ bằng}
  {$-5$}
  {$5$}{$0$}{$10$}

\cauTN[id=C4-B12-D01-NB-004,muc=NB,loigiai={$\displaystyle\int_{-1}^{3}\left[-3f(x)-4g(x)\right]\,\mathrm{d}x=-3\cdot \displaystyle\int_{-1}^{3}f(x)\,\mathrm{d}x-4\cdot \displaystyle\int_{-1}^{3}g(x)\,\mathrm{d}x=-10$.}]{Biết $\displaystyle\int_{-1}^{3}f(x)\,\mathrm{d}x=-6$ và $\displaystyle\int_{-1}^{3}g(x)\,\mathrm{d}x=7$. Giá trị của $\displaystyle\int_{-1}^{3}\left[-3f(x)-4g(x)\right]\,\mathrm{d}x$ bằng}
  {$-10$}
  {$46$}{$1$}{$3$}

\cauTN[id=C4-B12-D01-NB-005,muc=NB,loigiai={$\displaystyle\int_{-2}^{-1}\left[-4f(x)-3g(x)\right]\,\mathrm{d}x=-4\cdot \displaystyle\int_{-2}^{-1}f(x)\,\mathrm{d}x-3\cdot \displaystyle\int_{-2}^{-1}g(x)\,\mathrm{d}x=-48$.}]{Biết $\displaystyle\int_{-2}^{-1}f(x)\,\mathrm{d}x=9$ và $\displaystyle\int_{-2}^{-1}g(x)\,\mathrm{d}x=4$. Giá trị của $\displaystyle\int_{-2}^{-1}\left[-4f(x)-3g(x)\right]\,\mathrm{d}x$ bằng}
  {$-48$}
  {$-24$}{$13$}{$-43$}

\cauTN[id=C4-B12-D01-NB-006,muc=NB,loigiai={$\displaystyle\int_{-2}^{1}-5 f(x)\,\mathrm{d}x=-5\displaystyle\int_{-2}^{1}f(x)\,\mathrm{d}x=5$.}]{Biết $\displaystyle\int_{-2}^{1}f(x)\,\mathrm{d}x=-1$. Giá trị của $\displaystyle\int_{-2}^{1}-5 f(x)\,\mathrm{d}x$ bằng}
  {$5$}
  {$-1$}{$-6$}{$\dfrac{1}{5}$}

\cauTN[id=C4-B12-D01-NB-007,muc=NB,loigiai={$\displaystyle\int_{3}^{1}f(x)\,\mathrm{d}x=-\displaystyle\int_{1}^{3}f(x)\,\mathrm{d}x=-9$.}]{Biết $\displaystyle\int_{1}^{3}f(x)\,\mathrm{d}x=9$. Giá trị của $\displaystyle\int_{3}^{1}f(x)\,\mathrm{d}x$ bằng}
  {$-9$}
  {$9$}{$0$}{$18$}

\cauTN[id=C4-B12-D01-NB-008,muc=NB,loigiai={$\displaystyle\int_{3}^{0}f(x)\,\mathrm{d}x=-\displaystyle\int_{0}^{3}f(x)\,\mathrm{d}x=2$.}]{Biết $\displaystyle\int_{0}^{3}f(x)\,\mathrm{d}x=-2$. Giá trị của $\displaystyle\int_{3}^{0}f(x)\,\mathrm{d}x$ bằng}
  {$2$}
  {$-2$}{$0$}{$-4$}

\cauTN[id=C4-B12-D01-NB-009,muc=NB,loigiai={$\displaystyle\int_{2}^{4}\left[3f(x)+4g(x)\right]\,\mathrm{d}x=3\cdot \displaystyle\int_{2}^{4}f(x)\,\mathrm{d}x+4\cdot \displaystyle\int_{2}^{4}g(x)\,\mathrm{d}x=-34$.}]{Biết $\displaystyle\int_{2}^{4}f(x)\,\mathrm{d}x=-6$ và $\displaystyle\int_{2}^{4}g(x)\,\mathrm{d}x=-4$. Giá trị của $\displaystyle\int_{2}^{4}\left[3f(x)+4g(x)\right]\,\mathrm{d}x$ bằng}
  {$-34$}
  {$-2$}{$-10$}{$-36$}

\cauTN[id=C4-B12-D01-NB-010,muc=NB,loigiai={$\displaystyle\int_{1}^{0}f(x)\,\mathrm{d}x=-\displaystyle\int_{0}^{1}f(x)\,\mathrm{d}x=1$.}]{Biết $\displaystyle\int_{0}^{1}f(x)\,\mathrm{d}x=-1$. Giá trị của $\displaystyle\int_{1}^{0}f(x)\,\mathrm{d}x$ bằng}
  {$1$}
  {$-1$}{$0$}{$-2$}

\cauTN[id=C4-B12-D01-NB-011,muc=NB,loigiai={$\displaystyle\int_{-1}^{5}f(x)\,\mathrm{d}x=\displaystyle\int_{-1}^{2}f(x)\,\mathrm{d}x+\displaystyle\int_{2}^{5}f(x)\,\mathrm{d}x=5+(3)=8$.}]{Biết $\displaystyle\int_{-1}^{2}f(x)\,\mathrm{d}x=5$ và $\displaystyle\int_{2}^{5}f(x)\,\mathrm{d}x=3$. Giá trị của $\displaystyle\int_{-1}^{5}f(x)\,\mathrm{d}x$ bằng}
  {$8$}
  {$2$}{$-2$}{$15$}

\cauTN[id=C4-B12-D01-NB-012,muc=NB,loigiai={$\displaystyle\int_{1}^{5}\left[-3f(x)-g(x)\right]\,\mathrm{d}x=-3\cdot \displaystyle\int_{1}^{5}f(x)\,\mathrm{d}x-1\cdot \displaystyle\int_{1}^{5}g(x)\,\mathrm{d}x=-24$.}]{Biết $\displaystyle\int_{1}^{5}f(x)\,\mathrm{d}x=6$ và $\displaystyle\int_{1}^{5}g(x)\,\mathrm{d}x=6$. Giá trị của $\displaystyle\int_{1}^{5}\left[-3f(x)-g(x)\right]\,\mathrm{d}x$ bằng}
  {$-24$}
  {$-12$}{$12$}{$24$}

\cauTN[id=C4-B12-D01-NB-013,muc=NB,loigiai={$\displaystyle\int_{1}^{2}\left[f(x)-4g(x)\right]\,\mathrm{d}x=1\cdot \displaystyle\int_{1}^{2}f(x)\,\mathrm{d}x-4\cdot \displaystyle\int_{1}^{2}g(x)\,\mathrm{d}x=-2$.}]{Biết $\displaystyle\int_{1}^{2}f(x)\,\mathrm{d}x=2$ và $\displaystyle\int_{1}^{2}g(x)\,\mathrm{d}x=1$. Giá trị của $\displaystyle\int_{1}^{2}\left[f(x)-4g(x)\right]\,\mathrm{d}x$ bằng}
  {$-2$}
  {$6$}{$3$}{$-7$}

\cauTN[id=C4-B12-D01-NB-014,muc=NB,loigiai={$\displaystyle\int_{-2}^{1}\left[2f(x)+3g(x)\right]\,\mathrm{d}x=2\cdot \displaystyle\int_{-2}^{1}f(x)\,\mathrm{d}x+3\cdot \displaystyle\int_{-2}^{1}g(x)\,\mathrm{d}x=13$.}]{Biết $\displaystyle\int_{-2}^{1}f(x)\,\mathrm{d}x=-4$ và $\displaystyle\int_{-2}^{1}g(x)\,\mathrm{d}x=7$. Giá trị của $\displaystyle\int_{-2}^{1}\left[2f(x)+3g(x)\right]\,\mathrm{d}x$ bằng}
  {$13$}
  {$-29$}{$3$}{$2$}

\cauTN[id=C4-B12-D01-NB-015,muc=NB,loigiai={$\displaystyle\int_{2}^{0}f(x)\,\mathrm{d}x=-\displaystyle\int_{0}^{2}f(x)\,\mathrm{d}x=-3$.}]{Biết $\displaystyle\int_{0}^{2}f(x)\,\mathrm{d}x=3$. Giá trị của $\displaystyle\int_{2}^{0}f(x)\,\mathrm{d}x$ bằng}
  {$-3$}
  {$3$}{$0$}{$6$}

\cauTN[id=C4-B12-D01-NB-016,muc=NB,loigiai={$\displaystyle\int_{1}^{3}f(x)\,\mathrm{d}x=F(3)-F(1)=7-(-6)=13$.}]{Cho hàm số $f(x)$ liên tục trên $\mathbb{R}$ và $F(x)$ là một nguyên hàm của $f(x)$, thoả mãn $F(1)=-6$, $F(3)=7$. Giá trị của $\displaystyle\int_{1}^{3}f(x)\,\mathrm{d}x$ bằng}
  {$13$}
  {$-13$}{$1$}{$-42$}

\cauTN[id=C4-B12-D01-NB-017,muc=NB,loigiai={$\displaystyle\int_{1}^{-2}f(x)\,\mathrm{d}x=-\displaystyle\int_{-2}^{1}f(x)\,\mathrm{d}x=-5$.}]{Biết $\displaystyle\int_{-2}^{1}f(x)\,\mathrm{d}x=5$. Giá trị của $\displaystyle\int_{1}^{-2}f(x)\,\mathrm{d}x$ bằng}
  {$-5$}
  {$5$}{$0$}{$10$}

\cauTN[id=C4-B12-D01-NB-018,muc=NB,loigiai={$\displaystyle\int_{-1}^{1}f(x)\,\mathrm{d}x=F(1)-F(-1)=3-(-6)=9$.}]{Cho hàm số $f(x)$ liên tục trên $\mathbb{R}$ và $F(x)$ là một nguyên hàm của $f(x)$, thoả mãn $F(-1)=-6$, $F(1)=3$. Giá trị của $\displaystyle\int_{-1}^{1}f(x)\,\mathrm{d}x$ bằng}
  {$9$}
  {$-9$}{$-3$}{$-18$}

\cauTN[id=C4-B12-D01-NB-019,muc=NB,loigiai={$\displaystyle\int_{1}^{-2}f(x)\,\mathrm{d}x=-\displaystyle\int_{-2}^{1}f(x)\,\mathrm{d}x=2$.}]{Biết $\displaystyle\int_{-2}^{1}f(x)\,\mathrm{d}x=-2$. Giá trị của $\displaystyle\int_{1}^{-2}f(x)\,\mathrm{d}x$ bằng}
  {$2$}
  {$-2$}{$0$}{$-4$}

\cauTN[id=C4-B12-D01-NB-020,muc=NB,loigiai={$\displaystyle\int_{3}^{2}f(x)\,\mathrm{d}x=-\displaystyle\int_{2}^{3}f(x)\,\mathrm{d}x=5$.}]{Biết $\displaystyle\int_{2}^{3}f(x)\,\mathrm{d}x=-5$. Giá trị của $\displaystyle\int_{3}^{2}f(x)\,\mathrm{d}x$ bằng}
  {$5$}
  {$-5$}{$0$}{$-10$}

\cauTN[id=C4-B12-D01-NB-021,muc=NB,loigiai={$\displaystyle\int_{-1}^{1}2 f(x)\,\mathrm{d}x=2\displaystyle\int_{-1}^{1}f(x)\,\mathrm{d}x=16$.}]{Biết $\displaystyle\int_{-1}^{1}f(x)\,\mathrm{d}x=8$. Giá trị của $\displaystyle\int_{-1}^{1}2 f(x)\,\mathrm{d}x$ bằng}
  {$16$}
  {$8$}{$10$}{$4$}

\cauTN[id=C4-B12-D01-NB-022,muc=NB,loigiai={$\displaystyle\int_{2}^{3}f(x)\,\mathrm{d}x=F(3)-F(2)=5-(-2)=7$.}]{Cho hàm số $f(x)$ liên tục trên $\mathbb{R}$ và $F(x)$ là một nguyên hàm của $f(x)$, thoả mãn $F(2)=-2$, $F(3)=5$. Giá trị của $\displaystyle\int_{2}^{3}f(x)\,\mathrm{d}x$ bằng}
  {$7$}
  {$-7$}{$3$}{$-10$}

\cauTN[id=C4-B12-D01-NB-023,muc=NB,loigiai={$\displaystyle\int_{3}^{1}f(x)\,\mathrm{d}x=-\displaystyle\int_{1}^{3}f(x)\,\mathrm{d}x=7$.}]{Biết $\displaystyle\int_{1}^{3}f(x)\,\mathrm{d}x=-7$. Giá trị của $\displaystyle\int_{3}^{1}f(x)\,\mathrm{d}x$ bằng}
  {$7$}
  {$-7$}{$0$}{$-14$}

\cauTN[id=C4-B12-D01-NB-024,muc=NB,loigiai={$\displaystyle\int_{4}^{2}f(x)\,\mathrm{d}x=-\displaystyle\int_{2}^{4}f(x)\,\mathrm{d}x=3$.}]{Biết $\displaystyle\int_{2}^{4}f(x)\,\mathrm{d}x=-3$. Giá trị của $\displaystyle\int_{4}^{2}f(x)\,\mathrm{d}x$ bằng}
  {$3$}
  {$-3$}{$0$}{$-6$}

\cauTN[id=C4-B12-D01-NB-025,muc=NB,loigiai={$\displaystyle\int_{1}^{2}\left[-f(x)-g(x)\right]\,\mathrm{d}x=-1\cdot \displaystyle\int_{1}^{2}f(x)\,\mathrm{d}x-1\cdot \displaystyle\int_{1}^{2}g(x)\,\mathrm{d}x=2$.}]{Biết $\displaystyle\int_{1}^{2}f(x)\,\mathrm{d}x=-1$ và $\displaystyle\int_{1}^{2}g(x)\,\mathrm{d}x=-1$. Giá trị của $\displaystyle\int_{1}^{2}\left[-f(x)-g(x)\right]\,\mathrm{d}x$ bằng}
  {$2$}
  {$0$}{$-2$}{$-4$}

\cauTN[id=C4-B12-D01-NB-026,muc=NB,loigiai={$\displaystyle\int_{-2}^{2}f(x)\,\mathrm{d}x=F(2)-F(-2)=-4-(5)=-9$.}]{Cho hàm số $f(x)$ liên tục trên $\mathbb{R}$ và $F(x)$ là một nguyên hàm của $f(x)$, thoả mãn $F(-2)=5$, $F(2)=-4$. Giá trị của $\displaystyle\int_{-2}^{2}f(x)\,\mathrm{d}x$ bằng}
  {$-9$}
  {$9$}{$1$}{$-20$}

\cauTN[id=C4-B12-D01-NB-027,muc=NB,loigiai={$\displaystyle\int_{2}^{5}f(x)\,\mathrm{d}x=F(5)-F(2)=-4-(2)=-6$.}]{Cho hàm số $f(x)$ liên tục trên $\mathbb{R}$ và $F(x)$ là một nguyên hàm của $f(x)$, thoả mãn $F(2)=2$, $F(5)=-4$. Giá trị của $\displaystyle\int_{2}^{5}f(x)\,\mathrm{d}x$ bằng}
  {$-6$}
  {$6$}{$-2$}{$-8$}

\cauTN[id=C4-B12-D01-NB-028,muc=NB,loigiai={$\displaystyle\int_{1}^{4}-5 f(x)\,\mathrm{d}x=-5\displaystyle\int_{1}^{4}f(x)\,\mathrm{d}x=20$.}]{Biết $\displaystyle\int_{1}^{4}f(x)\,\mathrm{d}x=-4$. Giá trị của $\displaystyle\int_{1}^{4}-5 f(x)\,\mathrm{d}x$ bằng}
  {$20$}
  {$-4$}{$-9$}{$\dfrac{4}{5}$}

\cauTN[id=C4-B12-D01-NB-029,muc=NB,loigiai={$\displaystyle\int_{-1}^{-2}f(x)\,\mathrm{d}x=-\displaystyle\int_{-2}^{-1}f(x)\,\mathrm{d}x=2$.}]{Biết $\displaystyle\int_{-2}^{-1}f(x)\,\mathrm{d}x=-2$. Giá trị của $\displaystyle\int_{-1}^{-2}f(x)\,\mathrm{d}x$ bằng}
  {$2$}
  {$-2$}{$0$}{$-4$}

\cauTN[id=C4-B12-D01-NB-030,muc=NB,loigiai={$\displaystyle\int_{6}^{2}f(x)\,\mathrm{d}x=-\displaystyle\int_{2}^{6}f(x)\,\mathrm{d}x=-5$.}]{Biết $\displaystyle\int_{2}^{6}f(x)\,\mathrm{d}x=5$. Giá trị của $\displaystyle\int_{6}^{2}f(x)\,\mathrm{d}x$ bằng}
  {$-5$}
  {$5$}{$0$}{$10$}

\cauTN[id=C4-B12-D01-NB-031,muc=NB,loigiai={$\displaystyle\int_{2}^{4}\left[-4f(x)+3g(x)\right]\,\mathrm{d}x=-4\cdot \displaystyle\int_{2}^{4}f(x)\,\mathrm{d}x+3\cdot \displaystyle\int_{2}^{4}g(x)\,\mathrm{d}x=17$.}]{Biết $\displaystyle\int_{2}^{4}f(x)\,\mathrm{d}x=-2$ và $\displaystyle\int_{2}^{4}g(x)\,\mathrm{d}x=3$. Giá trị của $\displaystyle\int_{2}^{4}\left[-4f(x)+3g(x)\right]\,\mathrm{d}x$ bằng}
  {$17$}
  {$-1$}{$1$}{$-18$}

\cauTN[id=C4-B12-D01-NB-032,muc=NB,loigiai={$\displaystyle\int_{1}^{3}\left[-3f(x)+g(x)\right]\,\mathrm{d}x=-3\cdot \displaystyle\int_{1}^{3}f(x)\,\mathrm{d}x+1\cdot \displaystyle\int_{1}^{3}g(x)\,\mathrm{d}x=-20$.}]{Biết $\displaystyle\int_{1}^{3}f(x)\,\mathrm{d}x=8$ và $\displaystyle\int_{1}^{3}g(x)\,\mathrm{d}x=4$. Giá trị của $\displaystyle\int_{1}^{3}\left[-3f(x)+g(x)\right]\,\mathrm{d}x$ bằng}
  {$-20$}
  {$-28$}{$12$}{$-4$}

\cauTN[id=C4-B12-D01-NB-033,muc=NB,loigiai={$\displaystyle\int_{2}^{3}f(x)\,\mathrm{d}x=F(3)-F(2)=6-(7)=-1$.}]{Cho hàm số $f(x)$ liên tục trên $\mathbb{R}$ và $F(x)$ là một nguyên hàm của $f(x)$, thoả mãn $F(2)=7$, $F(3)=6$. Giá trị của $\displaystyle\int_{2}^{3}f(x)\,\mathrm{d}x$ bằng}
  {$-1$}
  {$1$}{$13$}{$42$}

\cauTN[id=C4-B12-D01-NB-034,muc=NB,loigiai={$\displaystyle\int_{-1}^{3}\left[-2f(x)+g(x)\right]\,\mathrm{d}x=-2\cdot \displaystyle\int_{-1}^{3}f(x)\,\mathrm{d}x+1\cdot \displaystyle\int_{-1}^{3}g(x)\,\mathrm{d}x=0$.}]{Biết $\displaystyle\int_{-1}^{3}f(x)\,\mathrm{d}x=-3$ và $\displaystyle\int_{-1}^{3}g(x)\,\mathrm{d}x=-6$. Giá trị của $\displaystyle\int_{-1}^{3}\left[-2f(x)+g(x)\right]\,\mathrm{d}x$ bằng}
  {$0$}
  {$12$}{$-9$}{$9$}

\cauTN[id=C4-B12-D01-NB-035,muc=NB,loigiai={$\displaystyle\int_{2}^{-2}f(x)\,\mathrm{d}x=-\displaystyle\int_{-2}^{2}f(x)\,\mathrm{d}x=-1$.}]{Biết $\displaystyle\int_{-2}^{2}f(x)\,\mathrm{d}x=1$. Giá trị của $\displaystyle\int_{2}^{-2}f(x)\,\mathrm{d}x$ bằng}
  {$-1$}
  {$1$}{$0$}{$2$}

\cauTN[id=C4-B12-D01-NB-036,muc=NB,loigiai={$\displaystyle\int_{-2}^{1}\left[3f(x)+3g(x)\right]\,\mathrm{d}x=3\cdot \displaystyle\int_{-2}^{1}f(x)\,\mathrm{d}x+3\cdot \displaystyle\int_{-2}^{1}g(x)\,\mathrm{d}x=15$.}]{Biết $\displaystyle\int_{-2}^{1}f(x)\,\mathrm{d}x=8$ và $\displaystyle\int_{-2}^{1}g(x)\,\mathrm{d}x=-3$. Giá trị của $\displaystyle\int_{-2}^{1}\left[3f(x)+3g(x)\right]\,\mathrm{d}x$ bằng}
  {$15$}
  {$33$}{$5$}{$-15$}

\cauDS[id=C4-B12-D01-NB-037,muc=NB,loigiai={Áp dụng các tính chất của tích phân. a) Sai: $\displaystyle\int_{0}^{4}f(x)\,\mathrm{d}x=3$. b) Sai: đúng là $6$. c) Đúng. d) Sai: đổi cận thì đổi dấu.}]{Cho hàm số $f(x)$ liên tục trên $\mathbb{R}$ thoả mãn $\displaystyle\int_{0}^{2}f(x)\,\mathrm{d}x=2$ và $\displaystyle\int_{2}^{4}f(x)\,\mathrm{d}x=1$. Xét tính đúng sai của các mệnh đề sau:}{%
  \y{$\displaystyle\int_{0}^{4}f(x)\,\mathrm{d}x=1$.}{s}
  \y{$\displaystyle\int_{0}^{2}3f(x)\,\mathrm{d}x=5$.}{s}
  \y{$\displaystyle\int_{0}^{4}f(x)\,\mathrm{d}x=3$.}{d}
  \y{$\displaystyle\int_{2}^{0}f(x)\,\mathrm{d}x=2$.}{s}}

\cauDS[id=C4-B12-D01-NB-038,muc=NB,loigiai={Áp dụng các tính chất của tích phân. a) Đúng. b) Đúng. c) Đúng. d) Sai: đúng là $15$.}]{Cho hàm số $f(x)$ liên tục trên $\mathbb{R}$ thoả mãn $\displaystyle\int_{-2}^{2}f(x)\,\mathrm{d}x=5$ và $\displaystyle\int_{2}^{3}f(x)\,\mathrm{d}x=3$. Xét tính đúng sai của các mệnh đề sau:}{%
  \y{$\displaystyle\int_{2}^{-2}f(x)\,\mathrm{d}x=-5$.}{d}
  \y{$\displaystyle\int_{-2}^{2}\left[f(x)+4\right]\,\mathrm{d}x=21$.}{d}
  \y{$\displaystyle\int_{-2}^{-2}f(x)\,\mathrm{d}x=0$.}{d}
  \y{$\displaystyle\int_{-2}^{2}3f(x)\,\mathrm{d}x=8$.}{s}}

\cauDS[id=C4-B12-D01-NB-039,muc=NB,loigiai={Áp dụng các tính chất của tích phân. a) Đúng. b) Sai: đổi cận thì đổi dấu. c) Sai: đúng là $7$. d) Sai: $\displaystyle\int_{2}^{6}f(x)\,\mathrm{d}x=4$.}]{Cho hàm số $f(x)$ liên tục trên $\mathbb{R}$ thoả mãn $\displaystyle\int_{2}^{3}f(x)\,\mathrm{d}x=9$ và $\displaystyle\int_{3}^{6}f(x)\,\mathrm{d}x=-5$. Xét tính đúng sai của các mệnh đề sau:}{%
  \y{$\displaystyle\int_{2}^{2}f(x)\,\mathrm{d}x=0$.}{d}
  \y{$\displaystyle\int_{3}^{2}f(x)\,\mathrm{d}x=9$.}{s}
  \y{$\displaystyle\int_{2}^{3}\left[f(x)-2\right]\,\mathrm{d}x=7$.}{s}
  \y{$\displaystyle\int_{2}^{6}f(x)\,\mathrm{d}x=14$.}{s}}

\cauDS[id=C4-B12-D01-NB-040,muc=NB,loigiai={Áp dụng các tính chất của tích phân. a) Sai: $\displaystyle\int_{0}^{4}f(x)\,\mathrm{d}x=8$. b) Đúng. c) Sai: đổi cận thì đổi dấu. d) Sai: đúng là $12$.}]{Cho hàm số $f(x)$ liên tục trên $\mathbb{R}$ thoả mãn $\displaystyle\int_{0}^{3}f(x)\,\mathrm{d}x=4$ và $\displaystyle\int_{3}^{4}f(x)\,\mathrm{d}x=4$. Xét tính đúng sai của các mệnh đề sau:}{%
  \y{$\displaystyle\int_{0}^{4}f(x)\,\mathrm{d}x=0$.}{s}
  \y{$\displaystyle\int_{0}^{3}3f(x)\,\mathrm{d}x=12$.}{d}
  \y{$\displaystyle\int_{3}^{0}f(x)\,\mathrm{d}x=4$.}{s}
  \y{$\displaystyle\int_{0}^{3}3f(x)\,\mathrm{d}x=7$.}{s}}

\cauDS[id=C4-B12-D01-NB-041,muc=NB,loigiai={Áp dụng các tính chất của tích phân. a) Sai: đúng là $-12$. b) Đúng. c) Đúng. d) Sai: đổi cận thì đổi dấu.}]{Cho hàm số $f(x)$ liên tục trên $\mathbb{R}$ thoả mãn $\displaystyle\int_{-1}^{2}f(x)\,\mathrm{d}x=-4$ và $\displaystyle\int_{2}^{3}f(x)\,\mathrm{d}x=8$. Xét tính đúng sai của các mệnh đề sau:}{%
  \y{$\displaystyle\int_{-1}^{2}3f(x)\,\mathrm{d}x=-1$.}{s}
  \y{$\displaystyle\int_{2}^{-1}f(x)\,\mathrm{d}x=4$.}{d}
  \y{$\displaystyle\int_{-1}^{2}3f(x)\,\mathrm{d}x=-12$.}{d}
  \y{$\displaystyle\int_{2}^{-1}f(x)\,\mathrm{d}x=-4$.}{s}}

\cauDS[id=C4-B12-D01-NB-042,muc=NB,loigiai={Áp dụng các tính chất của tích phân. a) Đúng. b) Sai: $\displaystyle\int_{-1}^{5}f(x)\,\mathrm{d}x=11$. c) Đúng. d) Đúng.}]{Cho hàm số $f(x)$ liên tục trên $\mathbb{R}$ thoả mãn $\displaystyle\int_{-1}^{2}f(x)\,\mathrm{d}x=3$ và $\displaystyle\int_{2}^{5}f(x)\,\mathrm{d}x=8$. Xét tính đúng sai của các mệnh đề sau:}{%
  \y{$\displaystyle\int_{-1}^{2}3f(x)\,\mathrm{d}x=9$.}{d}
  \y{$\displaystyle\int_{-1}^{5}f(x)\,\mathrm{d}x=-5$.}{s}
  \y{$\displaystyle\int_{-1}^{-1}f(x)\,\mathrm{d}x=0$.}{d}
  \y{$\displaystyle\int_{-1}^{2}\left[f(x)-5\right]\,\mathrm{d}x=-12$.}{d}}

\cauDS[id=C4-B12-D01-NB-043,muc=NB,loigiai={Áp dụng các tính chất của tích phân. a) Đúng. b) Sai: $\displaystyle\int_{1}^{5}f(x)\,\mathrm{d}x=12$. c) Sai: đúng là $3$. d) Đúng.}]{Cho hàm số $f(x)$ liên tục trên $\mathbb{R}$ thoả mãn $\displaystyle\int_{1}^{4}f(x)\,\mathrm{d}x=6$ và $\displaystyle\int_{4}^{5}f(x)\,\mathrm{d}x=6$. Xét tính đúng sai của các mệnh đề sau:}{%
  \y{$\displaystyle\int_{1}^{1}f(x)\,\mathrm{d}x=0$.}{d}
  \y{$\displaystyle\int_{1}^{5}f(x)\,\mathrm{d}x=0$.}{s}
  \y{$\displaystyle\int_{1}^{4}\left[f(x)-1\right]\,\mathrm{d}x=5$.}{s}
  \y{$\displaystyle\int_{4}^{1}f(x)\,\mathrm{d}x=-6$.}{d}}

\cauDS[id=C4-B12-D01-NB-044,muc=NB,loigiai={Áp dụng các tính chất của tích phân. a) Sai: đổi cận thì đổi dấu. b) Đúng. c) Đúng. d) Sai: $\displaystyle\int_{1}^{4}f(x)\,\mathrm{d}x=15$.}]{Cho hàm số $f(x)$ liên tục trên $\mathbb{R}$ thoả mãn $\displaystyle\int_{1}^{2}f(x)\,\mathrm{d}x=8$ và $\displaystyle\int_{2}^{4}f(x)\,\mathrm{d}x=7$. Xét tính đúng sai của các mệnh đề sau:}{%
  \y{$\displaystyle\int_{2}^{1}f(x)\,\mathrm{d}x=8$.}{s}
  \y{$\displaystyle\int_{2}^{1}f(x)\,\mathrm{d}x=-8$.}{d}
  \y{$\displaystyle\int_{1}^{2}\left[f(x)+2\right]\,\mathrm{d}x=10$.}{d}
  \y{$\displaystyle\int_{1}^{4}f(x)\,\mathrm{d}x=1$.}{s}}

\cauDS[id=C4-B12-D01-NB-045,muc=NB,loigiai={Áp dụng các tính chất của tích phân. a) Đúng. b) Sai: đúng là $3$. c) Đúng. d) Sai: $\displaystyle\int_{0}^{5}f(x)\,\mathrm{d}x=3$.}]{Cho hàm số $f(x)$ liên tục trên $\mathbb{R}$ thoả mãn $\displaystyle\int_{0}^{4}f(x)\,\mathrm{d}x=1$ và $\displaystyle\int_{4}^{5}f(x)\,\mathrm{d}x=2$. Xét tính đúng sai của các mệnh đề sau:}{%
  \y{$\displaystyle\int_{0}^{5}f(x)\,\mathrm{d}x=3$.}{d}
  \y{$\displaystyle\int_{0}^{4}3f(x)\,\mathrm{d}x=4$.}{s}
  \y{$\displaystyle\int_{0}^{4}3f(x)\,\mathrm{d}x=3$.}{d}
  \y{$\displaystyle\int_{0}^{5}f(x)\,\mathrm{d}x=-1$.}{s}}

\cauDS[id=C4-B12-D01-NB-046,muc=NB,loigiai={Áp dụng các tính chất của tích phân. a) Sai: $\displaystyle\int_{2}^{8}f(x)\,\mathrm{d}x=-2$. b) Đúng. c) Đúng. d) Sai: đổi cận thì đổi dấu.}]{Cho hàm số $f(x)$ liên tục trên $\mathbb{R}$ thoả mãn $\displaystyle\int_{2}^{6}f(x)\,\mathrm{d}x=1$ và $\displaystyle\int_{6}^{8}f(x)\,\mathrm{d}x=-3$. Xét tính đúng sai của các mệnh đề sau:}{%
  \y{$\displaystyle\int_{2}^{8}f(x)\,\mathrm{d}x=4$.}{s}
  \y{$\displaystyle\int_{2}^{8}f(x)\,\mathrm{d}x=-2$.}{d}
  \y{$\displaystyle\int_{2}^{2}f(x)\,\mathrm{d}x=0$.}{d}
  \y{$\displaystyle\int_{6}^{2}f(x)\,\mathrm{d}x=1$.}{s}}

\cauDS[id=C4-B12-D01-NB-047,muc=NB,loigiai={Áp dụng các tính chất của tích phân. a) Đúng. b) Sai: đúng là $24$. c) Sai: đổi cận thì đổi dấu. d) Sai: đúng là $9$.}]{Cho hàm số $f(x)$ liên tục trên $\mathbb{R}$ thoả mãn $\displaystyle\int_{1}^{2}f(x)\,\mathrm{d}x=8$ và $\displaystyle\int_{2}^{4}f(x)\,\mathrm{d}x=-4$. Xét tính đúng sai của các mệnh đề sau:}{%
  \y{$\displaystyle\int_{1}^{4}f(x)\,\mathrm{d}x=4$.}{d}
  \y{$\displaystyle\int_{1}^{2}3f(x)\,\mathrm{d}x=11$.}{s}
  \y{$\displaystyle\int_{2}^{1}f(x)\,\mathrm{d}x=8$.}{s}
  \y{$\displaystyle\int_{1}^{2}\left[f(x)+1\right]\,\mathrm{d}x=9$.}{s}}

\cauDS[id=C4-B12-D01-NB-048,muc=NB,loigiai={Áp dụng các tính chất của tích phân. a) Đúng. b) Đúng. c) Sai: $\displaystyle\int_{0}^{3}f(x)\,\mathrm{d}x=7$. d) Sai: đúng là $9$.}]{Cho hàm số $f(x)$ liên tục trên $\mathbb{R}$ thoả mãn $\displaystyle\int_{0}^{1}f(x)\,\mathrm{d}x=3$ và $\displaystyle\int_{1}^{3}f(x)\,\mathrm{d}x=4$. Xét tính đúng sai của các mệnh đề sau:}{%
  \y{$\displaystyle\int_{1}^{0}f(x)\,\mathrm{d}x=-3$.}{d}
  \y{$\displaystyle\int_{0}^{1}\left[f(x)+2\right]\,\mathrm{d}x=5$.}{d}
  \y{$\displaystyle\int_{0}^{3}f(x)\,\mathrm{d}x=-1$.}{s}
  \y{$\displaystyle\int_{0}^{1}3f(x)\,\mathrm{d}x=6$.}{s}}

\cauTLN[id=C4-B12-D01-NB-049,muc=NB,loigiai={$-1\cdot(-2)+(-2)\cdot(-3)=8$.}]{Biết $\displaystyle\int_{-1}^{1}f(x)\,\mathrm{d}x=-2$ và $\displaystyle\int_{-1}^{1}g(x)\,\mathrm{d}x=-3$. Tính $\displaystyle\int_{-1}^{1}\left[-f(x)-2g(x)\right]\,\mathrm{d}x$.}{8}

\cauTLN[id=C4-B12-D01-NB-050,muc=NB,loigiai={$\displaystyle\int_{0}^{3}f(x)\,\mathrm{d}x=\displaystyle\int_{-1}^{3}f(x)\,\mathrm{d}x-\displaystyle\int_{-1}^{0}f(x)\,\mathrm{d}x=4$.}]{Biết $\displaystyle\int_{-1}^{3}f(x)\,\mathrm{d}x=6$ và $\displaystyle\int_{-1}^{0}f(x)\,\mathrm{d}x=2$. Tính $\displaystyle\int_{0}^{3}f(x)\,\mathrm{d}x$.}{4}

\cauTLN[id=C4-B12-D01-NB-051,muc=NB,loigiai={$2\cdot(9)+(2)\cdot(7)=32$.}]{Biết $\displaystyle\int_{2}^{5}f(x)\,\mathrm{d}x=9$ và $\displaystyle\int_{2}^{5}g(x)\,\mathrm{d}x=7$. Tính $\displaystyle\int_{2}^{5}\left[2f(x)+2g(x)\right]\,\mathrm{d}x$.}{32}

\cauTLN[id=C4-B12-D01-NB-052,muc=NB,loigiai={$-2\cdot(5)+(3)\cdot(-6)=-28$.}]{Biết $\displaystyle\int_{2}^{6}f(x)\,\mathrm{d}x=5$ và $\displaystyle\int_{2}^{6}g(x)\,\mathrm{d}x=-6$. Tính $\displaystyle\int_{2}^{6}\left[-2f(x)+3g(x)\right]\,\mathrm{d}x$.}{-28}

\cauTLN[id=C4-B12-D01-NB-053,muc=NB,loigiai={$\displaystyle\int_{4}^{5}f(x)\,\mathrm{d}x=\displaystyle\int_{2}^{5}f(x)\,\mathrm{d}x-\displaystyle\int_{2}^{4}f(x)\,\mathrm{d}x=4$.}]{Biết $\displaystyle\int_{2}^{5}f(x)\,\mathrm{d}x=-3$ và $\displaystyle\int_{2}^{4}f(x)\,\mathrm{d}x=-7$. Tính $\displaystyle\int_{4}^{5}f(x)\,\mathrm{d}x$.}{4}

\cauTLN[id=C4-B12-D01-NB-054,muc=NB,loigiai={$4\cdot(1)+(-1)\cdot(1)=3$.}]{Biết $\displaystyle\int_{-2}^{0}f(x)\,\mathrm{d}x=1$ và $\displaystyle\int_{-2}^{0}g(x)\,\mathrm{d}x=1$. Tính $\displaystyle\int_{-2}^{0}\left[4f(x)-g(x)\right]\,\mathrm{d}x$.}{3}

\cauTLN[id=C4-B12-D01-NB-055,muc=NB,loigiai={$-2\cdot(-6)+(3)\cdot(-7)=-9$.}]{Biết $\displaystyle\int_{1}^{5}f(x)\,\mathrm{d}x=-6$ và $\displaystyle\int_{1}^{5}g(x)\,\mathrm{d}x=-7$. Tính $\displaystyle\int_{1}^{5}\left[-2f(x)+3g(x)\right]\,\mathrm{d}x$.}{-9}

\cauTLN[id=C4-B12-D01-NB-056,muc=NB,loigiai={$-2\cdot(-4)+(4)\cdot(1)=12$.}]{Biết $\displaystyle\int_{1}^{5}f(x)\,\mathrm{d}x=-4$ và $\displaystyle\int_{1}^{5}g(x)\,\mathrm{d}x=1$. Tính $\displaystyle\int_{1}^{5}\left[-2f(x)+4g(x)\right]\,\mathrm{d}x$.}{12}

\cauTLN[id=C4-B12-D01-NB-057,muc=NB,loigiai={$2\cdot(3)+(-3)\cdot(2)=0$.}]{Biết $\displaystyle\int_{2}^{6}f(x)\,\mathrm{d}x=3$ và $\displaystyle\int_{2}^{6}g(x)\,\mathrm{d}x=2$. Tính $\displaystyle\int_{2}^{6}\left[2f(x)-3g(x)\right]\,\mathrm{d}x$.}{0}

\cauTLN[id=C4-B12-D01-NB-058,muc=NB,loigiai={$-1\cdot(4)+(-3)\cdot(3)=-13$.}]{Biết $\displaystyle\int_{-1}^{3}f(x)\,\mathrm{d}x=4$ và $\displaystyle\int_{-1}^{3}g(x)\,\mathrm{d}x=3$. Tính $\displaystyle\int_{-1}^{3}\left[-f(x)-3g(x)\right]\,\mathrm{d}x$.}{-13}

\cauTLN[id=C4-B12-D01-NB-059,muc=NB,loigiai={$\displaystyle\int_{0}^{1}f(x)\,\mathrm{d}x=\displaystyle\int_{-2}^{1}f(x)\,\mathrm{d}x-\displaystyle\int_{-2}^{0}f(x)\,\mathrm{d}x=4$.}]{Biết $\displaystyle\int_{-2}^{1}f(x)\,\mathrm{d}x=2$ và $\displaystyle\int_{-2}^{0}f(x)\,\mathrm{d}x=-2$. Tính $\displaystyle\int_{0}^{1}f(x)\,\mathrm{d}x$.}{4}

\cauTLN[id=C4-B12-D01-NB-060,muc=NB,loigiai={$\displaystyle\int_{3}^{6}f(x)\,\mathrm{d}x=\displaystyle\int_{1}^{6}f(x)\,\mathrm{d}x-\displaystyle\int_{1}^{3}f(x)\,\mathrm{d}x=15$.}]{Biết $\displaystyle\int_{1}^{6}f(x)\,\mathrm{d}x=8$ và $\displaystyle\int_{1}^{3}f(x)\,\mathrm{d}x=-7$. Tính $\displaystyle\int_{3}^{6}f(x)\,\mathrm{d}x$.}{15}

\cauTN[id=C4-B12-D01-TH-001,muc=TH,loigiai={$\displaystyle\int_{0}^{1}f(x)\,\mathrm{d}x=\displaystyle\int_{-2}^{1}f(x)\,\mathrm{d}x-\displaystyle\int_{-2}^{0}f(x)\,\mathrm{d}x=3$ nên $\displaystyle\int_{0}^{1}-3 f(x)\,\mathrm{d}x=-9$.}]{Cho hàm số $f(x)$ liên tục trên $\mathbb{R}$, $\displaystyle\int_{-2}^{0}f(x)\,\mathrm{d}x=2$ và $\displaystyle\int_{-2}^{1}f(x)\,\mathrm{d}x=5$. Giá trị của $\displaystyle\int_{0}^{1}-3 f(x)\,\mathrm{d}x$ bằng}
  {$-9$}
  {$-21$}{$9$}{$3$}

\cauTN[id=C4-B12-D01-TH-002,muc=TH,loigiai={$\displaystyle\int_{6}^{7}f(x)\,\mathrm{d}x=\displaystyle\int_{2}^{7}f(x)\,\mathrm{d}x-\displaystyle\int_{2}^{6}f(x)\,\mathrm{d}x=11$ nên $\displaystyle\int_{6}^{7}3 f(x)\,\mathrm{d}x=33$.}]{Cho hàm số $f(x)$ liên tục trên $\mathbb{R}$, $\displaystyle\int_{2}^{6}f(x)\,\mathrm{d}x=-5$ và $\displaystyle\int_{2}^{7}f(x)\,\mathrm{d}x=6$. Giá trị của $\displaystyle\int_{6}^{7}3 f(x)\,\mathrm{d}x$ bằng}
  {$33$}
  {$3$}{$-33$}{$11$}

\cauTN[id=C4-B12-D01-TH-003,muc=TH,loigiai={$\displaystyle\int_{2}^{3}f(x)\,\mathrm{d}x=\displaystyle\int_{-2}^{3}f(x)\,\mathrm{d}x-\displaystyle\int_{-2}^{2}f(x)\,\mathrm{d}x=-1$ nên $\displaystyle\int_{2}^{3}-1 f(x)\,\mathrm{d}x=1$.}]{Cho hàm số $f(x)$ liên tục trên $\mathbb{R}$, $\displaystyle\int_{-2}^{2}f(x)\,\mathrm{d}x=-3$ và $\displaystyle\int_{-2}^{3}f(x)\,\mathrm{d}x=-4$. Giá trị của $\displaystyle\int_{2}^{3}-1 f(x)\,\mathrm{d}x$ bằng}
  {$1$}
  {$7$}{$-1$}{$4$}

\cauTN[id=C4-B12-D01-TH-004,muc=TH,loigiai={$\displaystyle\int_{2}^{4}f'(x)\,\mathrm{d}x=f(4)-f(2)=-5$.}]{Cho hàm số $f(x)$ có đạo hàm liên tục trên $\mathbb{R}$, $f(2)=2$ và $f(4)=-3$. Giá trị của $\displaystyle\int_{2}^{4}f'(x)\,\mathrm{d}x$ bằng}
  {$-5$}
  {$5$}{$-1$}{$-3$}

\cauTN[id=C4-B12-D01-TH-005,muc=TH,loigiai={$\displaystyle\int_{0}^{3}f(x)\,\mathrm{d}x=\displaystyle\int_{-2}^{3}f(x)\,\mathrm{d}x-\displaystyle\int_{-2}^{0}f(x)\,\mathrm{d}x=-14$ nên $\displaystyle\int_{0}^{3}-3 f(x)\,\mathrm{d}x=42$.}]{Cho hàm số $f(x)$ liên tục trên $\mathbb{R}$, $\displaystyle\int_{-2}^{0}f(x)\,\mathrm{d}x=7$ và $\displaystyle\int_{-2}^{3}f(x)\,\mathrm{d}x=-7$. Giá trị của $\displaystyle\int_{0}^{3}-3 f(x)\,\mathrm{d}x$ bằng}
  {$42$}
  {$0$}{$-42$}{$-14$}

\cauTN[id=C4-B12-D01-TH-006,muc=TH,loigiai={$\displaystyle\int_{-1}^{2}\left[-2f(x)-3x-3\right]\,\mathrm{d}x=-2\displaystyle\int_{-1}^{2}f(x)\,\mathrm{d}x+\displaystyle\int_{-1}^{2}(-3x-3)\,\mathrm{d}x=-2\cdot(-5)-\dfrac{27}{2}=-\dfrac{7}{2}$.}]{Biết $\displaystyle\int_{-1}^{2}f(x)\,\mathrm{d}x=-5$. Giá trị của $\displaystyle\int_{-1}^{2}\left[-2f(x)-3x-3\right]\,\mathrm{d}x$ bằng}
  {$-\dfrac{7}{2}$}
  {$-8$}{$4$}{$-\dfrac{37}{2}$}

\cauTN[id=C4-B12-D01-TH-007,muc=TH,loigiai={$\displaystyle\int_{-1}^{3}f'(x)\,\mathrm{d}x=f(3)-f(-1)=1$.}]{Cho hàm số $f(x)$ có đạo hàm liên tục trên $\mathbb{R}$, $f(-1)=4$ và $f(3)=5$. Giá trị của $\displaystyle\int_{-1}^{3}f'(x)\,\mathrm{d}x$ bằng}
  {$1$}
  {$-1$}{$9$}{$5$}

\cauTN[id=C4-B12-D01-TH-008,muc=TH,loigiai={$\displaystyle\int_{3}^{4}f(x)\,\mathrm{d}x=\displaystyle\int_{-1}^{4}f(x)\,\mathrm{d}x-\displaystyle\int_{-1}^{3}f(x)\,\mathrm{d}x=1$ nên $\displaystyle\int_{3}^{4}-3 f(x)\,\mathrm{d}x=-3$.}]{Cho hàm số $f(x)$ liên tục trên $\mathbb{R}$, $\displaystyle\int_{-1}^{3}f(x)\,\mathrm{d}x=-4$ và $\displaystyle\int_{-1}^{4}f(x)\,\mathrm{d}x=-3$. Giá trị của $\displaystyle\int_{3}^{4}-3 f(x)\,\mathrm{d}x$ bằng}
  {$-3$}
  {$21$}{$3$}{$1$}

\cauTN[id=C4-B12-D01-TH-009,muc=TH,loigiai={$\displaystyle\int_{-2}^{0}f'(x)\,\mathrm{d}x=f(0)-f(-2)=4$.}]{Cho hàm số $f(x)$ có đạo hàm liên tục trên $\mathbb{R}$, $f(-2)=3$ và $f(0)=7$. Giá trị của $\displaystyle\int_{-2}^{0}f'(x)\,\mathrm{d}x$ bằng}
  {$4$}
  {$-4$}{$10$}{$7$}

\cauTN[id=C4-B12-D01-TH-010,muc=TH,loigiai={$\displaystyle\int_{-2}^{-1}f'(x)\,\mathrm{d}x=f(-1)-f(-2)=-3$.}]{Cho hàm số $f(x)$ có đạo hàm liên tục trên $\mathbb{R}$, $f(-2)=3$ và $f(-1)=0$. Giá trị của $\displaystyle\int_{-2}^{-1}f'(x)\,\mathrm{d}x$ bằng}
  {$-3$}
  {$3$}{$0$}{$-6$}

\cauTN[id=C4-B12-D01-TH-011,muc=TH,loigiai={$\displaystyle\int_{1}^{2}f'(x)\,\mathrm{d}x=f(2)-f(1)=5$.}]{Cho hàm số $f(x)$ có đạo hàm liên tục trên $\mathbb{R}$, $f(1)=1$ và $f(2)=6$. Giá trị của $\displaystyle\int_{1}^{2}f'(x)\,\mathrm{d}x$ bằng}
  {$5$}
  {$-5$}{$7$}{$6$}

\cauTN[id=C4-B12-D01-TH-012,muc=TH,loigiai={$\displaystyle\int_{1}^{3}f(x)\,\mathrm{d}x=\displaystyle\int_{-2}^{3}f(x)\,\mathrm{d}x-\displaystyle\int_{-2}^{1}f(x)\,\mathrm{d}x=12$ nên $\displaystyle\int_{1}^{3}2 f(x)\,\mathrm{d}x=24$.}]{Cho hàm số $f(x)$ liên tục trên $\mathbb{R}$, $\displaystyle\int_{-2}^{1}f(x)\,\mathrm{d}x=-4$ và $\displaystyle\int_{-2}^{3}f(x)\,\mathrm{d}x=8$. Giá trị của $\displaystyle\int_{1}^{3}2 f(x)\,\mathrm{d}x$ bằng}
  {$24$}
  {$8$}{$-24$}{$12$}

\cauTN[id=C4-B12-D01-TH-013,muc=TH,loigiai={$\displaystyle\int_{2}^{6}f'(x)\,\mathrm{d}x=f(6)-f(2)=-5$.}]{Cho hàm số $f(x)$ có đạo hàm liên tục trên $\mathbb{R}$, $f(2)=2$ và $f(6)=-3$. Giá trị của $\displaystyle\int_{2}^{6}f'(x)\,\mathrm{d}x$ bằng}
  {$-5$}
  {$5$}{$-1$}{$-3$}

\cauTN[id=C4-B12-D01-TH-014,muc=TH,loigiai={$\displaystyle\int_{-2}^{-1}\left[2f(x)-x+3\right]\,\mathrm{d}x=2\displaystyle\int_{-2}^{-1}f(x)\,\mathrm{d}x+\displaystyle\int_{-2}^{-1}(-x+3)\,\mathrm{d}x=2\cdot(-4)+\dfrac{9}{2}=-\dfrac{7}{2}$.}]{Biết $\displaystyle\int_{-2}^{-1}f(x)\,\mathrm{d}x=-4$. Giá trị của $\displaystyle\int_{-2}^{-1}\left[2f(x)-x+3\right]\,\mathrm{d}x$ bằng}
  {$-\dfrac{7}{2}$}
  {$-2$}{$-6$}{$\dfrac{1}{2}$}

\cauTN[id=C4-B12-D01-TH-015,muc=TH,loigiai={$\displaystyle\int_{-1}^{0}\left[-3f(x)-2x-1\right]\,\mathrm{d}x=-3\displaystyle\int_{-1}^{0}f(x)\,\mathrm{d}x+\displaystyle\int_{-1}^{0}(-2x-1)\,\mathrm{d}x=-3\cdot(7)+0=-21$.}]{Biết $\displaystyle\int_{-1}^{0}f(x)\,\mathrm{d}x=7$. Giá trị của $\displaystyle\int_{-1}^{0}\left[-3f(x)-2x-1\right]\,\mathrm{d}x$ bằng}
  {$-21$}
  {$-20$}{$-24$}{$7$}

\cauTN[id=C4-B12-D01-TH-016,muc=TH,loigiai={$\displaystyle\int_{1}^{4}f'(x)\,\mathrm{d}x=f(4)-f(1)=1$.}]{Cho hàm số $f(x)$ có đạo hàm liên tục trên $\mathbb{R}$, $f(1)=-4$ và $f(4)=-3$. Giá trị của $\displaystyle\int_{1}^{4}f'(x)\,\mathrm{d}x$ bằng}
  {$1$}
  {$-1$}{$-7$}{$-3$}

\cauTN[id=C4-B12-D01-TH-017,muc=TH,loigiai={$\displaystyle\int_{-2}^{1}f'(x)\,\mathrm{d}x=f(1)-f(-2)=-2$.}]{Cho hàm số $f(x)$ có đạo hàm liên tục trên $\mathbb{R}$, $f(-2)=-2$ và $f(1)=-4$. Giá trị của $\displaystyle\int_{-2}^{1}f'(x)\,\mathrm{d}x$ bằng}
  {$-2$}
  {$2$}{$-6$}{$-4$}

\cauTN[id=C4-B12-D01-TH-018,muc=TH,loigiai={$\displaystyle\int_{2}^{5}f(x)\,\mathrm{d}x=\displaystyle\int_{1}^{5}f(x)\,\mathrm{d}x-\displaystyle\int_{1}^{2}f(x)\,\mathrm{d}x=-15$ nên $\displaystyle\int_{2}^{5}4 f(x)\,\mathrm{d}x=-60$.}]{Cho hàm số $f(x)$ liên tục trên $\mathbb{R}$, $\displaystyle\int_{1}^{2}f(x)\,\mathrm{d}x=9$ và $\displaystyle\int_{1}^{5}f(x)\,\mathrm{d}x=-6$. Giá trị của $\displaystyle\int_{2}^{5}4 f(x)\,\mathrm{d}x$ bằng}
  {$-60$}
  {$12$}{$60$}{$-15$}

\cauTN[id=C4-B12-D01-TH-019,muc=TH,loigiai={$\displaystyle\int_{2}^{3}f(x)\,\mathrm{d}x=\displaystyle\int_{-2}^{3}f(x)\,\mathrm{d}x-\displaystyle\int_{-2}^{2}f(x)\,\mathrm{d}x=9$ nên $\displaystyle\int_{2}^{3}4 f(x)\,\mathrm{d}x=36$.}]{Cho hàm số $f(x)$ liên tục trên $\mathbb{R}$, $\displaystyle\int_{-2}^{2}f(x)\,\mathrm{d}x=-2$ và $\displaystyle\int_{-2}^{3}f(x)\,\mathrm{d}x=7$. Giá trị của $\displaystyle\int_{2}^{3}4 f(x)\,\mathrm{d}x$ bằng}
  {$36$}
  {$20$}{$-36$}{$9$}

\cauTN[id=C4-B12-D01-TH-020,muc=TH,loigiai={$\displaystyle\int_{2}^{4}f(x)\,\mathrm{d}x=\displaystyle\int_{1}^{4}f(x)\,\mathrm{d}x-\displaystyle\int_{1}^{2}f(x)\,\mathrm{d}x=1$ nên $\displaystyle\int_{2}^{4}-3 f(x)\,\mathrm{d}x=-3$.}]{Cho hàm số $f(x)$ liên tục trên $\mathbb{R}$, $\displaystyle\int_{1}^{2}f(x)\,\mathrm{d}x=4$ và $\displaystyle\int_{1}^{4}f(x)\,\mathrm{d}x=5$. Giá trị của $\displaystyle\int_{2}^{4}-3 f(x)\,\mathrm{d}x$ bằng}
  {$-3$}
  {$-27$}{$3$}{$1$}

\cauTN[id=C4-B12-D01-TH-021,muc=TH,loigiai={$\displaystyle\int_{1}^{2}f'(x)\,\mathrm{d}x=f(2)-f(1)=1$.}]{Cho hàm số $f(x)$ có đạo hàm liên tục trên $\mathbb{R}$, $f(1)=-4$ và $f(2)=-3$. Giá trị của $\displaystyle\int_{1}^{2}f'(x)\,\mathrm{d}x$ bằng}
  {$1$}
  {$-1$}{$-7$}{$-3$}

\cauTN[id=C4-B12-D01-TH-022,muc=TH,loigiai={$\displaystyle\int_{0}^{2}\left[2f(x)+x\right]\,\mathrm{d}x=2\displaystyle\int_{0}^{2}f(x)\,\mathrm{d}x+\displaystyle\int_{0}^{2}(x)\,\mathrm{d}x=2\cdot(-3)+2=-4$.}]{Biết $\displaystyle\int_{0}^{2}f(x)\,\mathrm{d}x=-3$. Giá trị của $\displaystyle\int_{0}^{2}\left[2f(x)+x\right]\,\mathrm{d}x$ bằng}
  {$-4$}
  {$-2$}{$-5$}{$-1$}

\cauTN[id=C4-B12-D01-TH-023,muc=TH,loigiai={$\displaystyle\int_{2}^{3}f(x)\,\mathrm{d}x=\displaystyle\int_{-1}^{3}f(x)\,\mathrm{d}x-\displaystyle\int_{-1}^{2}f(x)\,\mathrm{d}x=3$ nên $\displaystyle\int_{2}^{3}-2 f(x)\,\mathrm{d}x=-6$.}]{Cho hàm số $f(x)$ liên tục trên $\mathbb{R}$, $\displaystyle\int_{-1}^{2}f(x)\,\mathrm{d}x=5$ và $\displaystyle\int_{-1}^{3}f(x)\,\mathrm{d}x=8$. Giá trị của $\displaystyle\int_{2}^{3}-2 f(x)\,\mathrm{d}x$ bằng}
  {$-6$}
  {$-26$}{$6$}{$3$}

\cauTN[id=C4-B12-D01-TH-024,muc=TH,loigiai={$\displaystyle\int_{2}^{6}f'(x)\,\mathrm{d}x=f(6)-f(2)=2$.}]{Cho hàm số $f(x)$ có đạo hàm liên tục trên $\mathbb{R}$, $f(2)=3$ và $f(6)=5$. Giá trị của $\displaystyle\int_{2}^{6}f'(x)\,\mathrm{d}x$ bằng}
  {$2$}
  {$-2$}{$8$}{$5$}

\cauTN[id=C4-B12-D01-TH-025,muc=TH,loigiai={$\displaystyle\int_{0}^{4}\left[3f(x)-x+2\right]\,\mathrm{d}x=3\displaystyle\int_{0}^{4}f(x)\,\mathrm{d}x+\displaystyle\int_{0}^{4}(-x+2)\,\mathrm{d}x=3\cdot(1)+0=3$.}]{Biết $\displaystyle\int_{0}^{4}f(x)\,\mathrm{d}x=1$. Giá trị của $\displaystyle\int_{0}^{4}\left[3f(x)-x+2\right]\,\mathrm{d}x$ bằng}
  {$3$}
  {$-5$}{$4$}{$1$}

\cauTN[id=C4-B12-D01-TH-026,muc=TH,loigiai={$\displaystyle\int_{1}^{3}\left[3f(x)-2x-3\right]\,\mathrm{d}x=3\displaystyle\int_{1}^{3}f(x)\,\mathrm{d}x+\displaystyle\int_{1}^{3}(-2x-3)\,\mathrm{d}x=3\cdot(4)-14=-2$.}]{Biết $\displaystyle\int_{1}^{3}f(x)\,\mathrm{d}x=4$. Giá trị của $\displaystyle\int_{1}^{3}\left[3f(x)-2x-3\right]\,\mathrm{d}x$ bằng}
  {$-2$}
  {$-10$}{$7$}{$1$}

\cauTN[id=C4-B12-D01-TH-027,muc=TH,loigiai={$\displaystyle\int_{3}^{4}f(x)\,\mathrm{d}x=\displaystyle\int_{0}^{4}f(x)\,\mathrm{d}x-\displaystyle\int_{0}^{3}f(x)\,\mathrm{d}x=-10$ nên $\displaystyle\int_{3}^{4}3 f(x)\,\mathrm{d}x=-30$.}]{Cho hàm số $f(x)$ liên tục trên $\mathbb{R}$, $\displaystyle\int_{0}^{3}f(x)\,\mathrm{d}x=8$ và $\displaystyle\int_{0}^{4}f(x)\,\mathrm{d}x=-2$. Giá trị của $\displaystyle\int_{3}^{4}3 f(x)\,\mathrm{d}x$ bằng}
  {$-30$}
  {$18$}{$30$}{$-10$}

\cauTN[id=C4-B12-D01-TH-028,muc=TH,loigiai={$\displaystyle\int_{3}^{5}f(x)\,\mathrm{d}x=\displaystyle\int_{1}^{5}f(x)\,\mathrm{d}x-\displaystyle\int_{1}^{3}f(x)\,\mathrm{d}x=-2$ nên $\displaystyle\int_{3}^{5}-1 f(x)\,\mathrm{d}x=2$.}]{Cho hàm số $f(x)$ liên tục trên $\mathbb{R}$, $\displaystyle\int_{1}^{3}f(x)\,\mathrm{d}x=-2$ và $\displaystyle\int_{1}^{5}f(x)\,\mathrm{d}x=-4$. Giá trị của $\displaystyle\int_{3}^{5}-1 f(x)\,\mathrm{d}x$ bằng}
  {$2$}
  {$6$}{$-2$}{$4$}

\cauTN[id=C4-B12-D01-TH-029,muc=TH,loigiai={$\displaystyle\int_{-1}^{2}f'(x)\,\mathrm{d}x=f(2)-f(-1)=-1$.}]{Cho hàm số $f(x)$ có đạo hàm liên tục trên $\mathbb{R}$, $f(-1)=0$ và $f(2)=-1$. Giá trị của $\displaystyle\int_{-1}^{2}f'(x)\,\mathrm{d}x$ bằng}
  {$-1$}
  {$1$}{$0$}{$-2$}

\cauTN[id=C4-B12-D01-TH-030,muc=TH,loigiai={$\displaystyle\int_{1}^{3}f'(x)\,\mathrm{d}x=f(3)-f(1)=-5$.}]{Cho hàm số $f(x)$ có đạo hàm liên tục trên $\mathbb{R}$, $f(1)=4$ và $f(3)=-1$. Giá trị của $\displaystyle\int_{1}^{3}f'(x)\,\mathrm{d}x$ bằng}
  {$-5$}
  {$5$}{$3$}{$-1$}

\cauTN[id=C4-B12-D01-TH-031,muc=TH,loigiai={$\displaystyle\int_{1}^{5}\left[-2f(x)-3x\right]\,\mathrm{d}x=-2\displaystyle\int_{1}^{5}f(x)\,\mathrm{d}x+\displaystyle\int_{1}^{5}(-3x)\,\mathrm{d}x=-2\cdot(6)-36=-48$.}]{Biết $\displaystyle\int_{1}^{5}f(x)\,\mathrm{d}x=6$. Giá trị của $\displaystyle\int_{1}^{5}\left[-2f(x)-3x\right]\,\mathrm{d}x$ bằng}
  {$-48$}
  {$-84$}{$-15$}{$-30$}

\cauTN[id=C4-B12-D01-TH-032,muc=TH,loigiai={$\displaystyle\int_{-1}^{1}f'(x)\,\mathrm{d}x=f(1)-f(-1)=5$.}]{Cho hàm số $f(x)$ có đạo hàm liên tục trên $\mathbb{R}$, $f(-1)=2$ và $f(1)=7$. Giá trị của $\displaystyle\int_{-1}^{1}f'(x)\,\mathrm{d}x$ bằng}
  {$5$}
  {$-5$}{$9$}{$7$}

\cauTN[id=C4-B12-D01-TH-033,muc=TH,loigiai={$\displaystyle\int_{2}^{4}f(x)\,\mathrm{d}x=\displaystyle\int_{0}^{4}f(x)\,\mathrm{d}x-\displaystyle\int_{0}^{2}f(x)\,\mathrm{d}x=-6$ nên $\displaystyle\int_{2}^{4}3 f(x)\,\mathrm{d}x=-18$.}]{Cho hàm số $f(x)$ liên tục trên $\mathbb{R}$, $\displaystyle\int_{0}^{2}f(x)\,\mathrm{d}x=4$ và $\displaystyle\int_{0}^{4}f(x)\,\mathrm{d}x=-2$. Giá trị của $\displaystyle\int_{2}^{4}3 f(x)\,\mathrm{d}x$ bằng}
  {$-18$}
  {$6$}{$18$}{$-6$}

\cauTN[id=C4-B12-D01-TH-034,muc=TH,loigiai={$\displaystyle\int_{1}^{3}\left[2f(x)-2x+1\right]\,\mathrm{d}x=2\displaystyle\int_{1}^{3}f(x)\,\mathrm{d}x+\displaystyle\int_{1}^{3}(-2x+1)\,\mathrm{d}x=2\cdot(2)-6=-2$.}]{Biết $\displaystyle\int_{1}^{3}f(x)\,\mathrm{d}x=2$. Giá trị của $\displaystyle\int_{1}^{3}\left[2f(x)-2x+1\right]\,\mathrm{d}x$ bằng}
  {$-2$}
  {$-10$}{$3$}{$-4$}

\cauTN[id=C4-B12-D01-TH-035,muc=TH,loigiai={$\displaystyle\int_{-2}^{1}\left[-2f(x)-x-3\right]\,\mathrm{d}x=-2\displaystyle\int_{-2}^{1}f(x)\,\mathrm{d}x+\displaystyle\int_{-2}^{1}(-x-3)\,\mathrm{d}x=-2\cdot(-3)-\dfrac{15}{2}=-\dfrac{3}{2}$.}]{Biết $\displaystyle\int_{-2}^{1}f(x)\,\mathrm{d}x=-3$. Giá trị của $\displaystyle\int_{-2}^{1}\left[-2f(x)-x-3\right]\,\mathrm{d}x$ bằng}
  {$-\dfrac{3}{2}$}
  {$0$}{$2$}{$-\dfrac{21}{2}$}

\cauTN[id=C4-B12-D01-TH-036,muc=TH,loigiai={$\displaystyle\int_{0}^{3}f'(x)\,\mathrm{d}x=f(3)-f(0)=-2$.}]{Cho hàm số $f(x)$ có đạo hàm liên tục trên $\mathbb{R}$, $f(0)=1$ và $f(3)=-1$. Giá trị của $\displaystyle\int_{0}^{3}f'(x)\,\mathrm{d}x$ bằng}
  {$-2$}
  {$2$}{$0$}{$-1$}

\cauTN[id=C4-B12-D01-TH-037,muc=TH,loigiai={$\displaystyle\int_{-2}^{2}\left[4f(x)+3x-4\right]\,\mathrm{d}x=4\displaystyle\int_{-2}^{2}f(x)\,\mathrm{d}x+\displaystyle\int_{-2}^{2}(3x-4)\,\mathrm{d}x=4\cdot(2)-16=-8$.}]{Biết $\displaystyle\int_{-2}^{2}f(x)\,\mathrm{d}x=2$. Giá trị của $\displaystyle\int_{-2}^{2}\left[4f(x)+3x-4\right]\,\mathrm{d}x$ bằng}
  {$-8$}
  {$7$}{$-14$}{$4$}

\cauTN[id=C4-B12-D01-TH-038,muc=TH,loigiai={$\displaystyle\int_{0}^{2}f'(x)\,\mathrm{d}x=f(2)-f(0)=9$.}]{Cho hàm số $f(x)$ có đạo hàm liên tục trên $\mathbb{R}$, $f(0)=-4$ và $f(2)=5$. Giá trị của $\displaystyle\int_{0}^{2}f'(x)\,\mathrm{d}x$ bằng}
  {$9$}
  {$-9$}{$1$}{$5$}

\cauTN[id=C4-B12-D01-TH-039,muc=TH,loigiai={$\displaystyle\int_{0}^{1}f'(x)\,\mathrm{d}x=f(1)-f(0)=3$.}]{Cho hàm số $f(x)$ có đạo hàm liên tục trên $\mathbb{R}$, $f(0)=4$ và $f(1)=7$. Giá trị của $\displaystyle\int_{0}^{1}f'(x)\,\mathrm{d}x$ bằng}
  {$3$}
  {$-3$}{$11$}{$7$}

\cauTN[id=C4-B12-D01-TH-040,muc=TH,loigiai={$\displaystyle\int_{0}^{1}f(x)\,\mathrm{d}x=\displaystyle\int_{-1}^{1}f(x)\,\mathrm{d}x-\displaystyle\int_{-1}^{0}f(x)\,\mathrm{d}x=7$ nên $\displaystyle\int_{0}^{1}3 f(x)\,\mathrm{d}x=21$.}]{Cho hàm số $f(x)$ liên tục trên $\mathbb{R}$, $\displaystyle\int_{-1}^{0}f(x)\,\mathrm{d}x=-2$ và $\displaystyle\int_{-1}^{1}f(x)\,\mathrm{d}x=5$. Giá trị của $\displaystyle\int_{0}^{1}3 f(x)\,\mathrm{d}x$ bằng}
  {$21$}
  {$9$}{$-21$}{$7$}

\cauTN[id=C4-B12-D01-TH-041,muc=TH,loigiai={$\displaystyle\int_{0}^{2}f'(x)\,\mathrm{d}x=f(2)-f(0)=5$.}]{Cho hàm số $f(x)$ có đạo hàm liên tục trên $\mathbb{R}$, $f(0)=2$ và $f(2)=7$. Giá trị của $\displaystyle\int_{0}^{2}f'(x)\,\mathrm{d}x$ bằng}
  {$5$}
  {$-5$}{$9$}{$7$}

\cauTN[id=C4-B12-D01-TH-042,muc=TH,loigiai={$\displaystyle\int_{-1}^{0}\left[-f(x)+x+4\right]\,\mathrm{d}x=-1\displaystyle\int_{-1}^{0}f(x)\,\mathrm{d}x+\displaystyle\int_{-1}^{0}(x+4)\,\mathrm{d}x=-1\cdot(8)+\dfrac{7}{2}=-\dfrac{9}{2}$.}]{Biết $\displaystyle\int_{-1}^{0}f(x)\,\mathrm{d}x=8$. Giá trị của $\displaystyle\int_{-1}^{0}\left[-f(x)+x+4\right]\,\mathrm{d}x$ bằng}
  {$-\dfrac{9}{2}$}
  {$-5$}{$-3$}{$\dfrac{23}{2}$}

\cauTN[id=C4-B12-D01-TH-043,muc=TH,loigiai={$\displaystyle\int_{5}^{7}f(x)\,\mathrm{d}x=\displaystyle\int_{1}^{7}f(x)\,\mathrm{d}x-\displaystyle\int_{1}^{5}f(x)\,\mathrm{d}x=-9$ nên $\displaystyle\int_{5}^{7}2 f(x)\,\mathrm{d}x=-18$.}]{Cho hàm số $f(x)$ liên tục trên $\mathbb{R}$, $\displaystyle\int_{1}^{5}f(x)\,\mathrm{d}x=8$ và $\displaystyle\int_{1}^{7}f(x)\,\mathrm{d}x=-1$. Giá trị của $\displaystyle\int_{5}^{7}2 f(x)\,\mathrm{d}x$ bằng}
  {$-18$}
  {$14$}{$18$}{$-9$}

\cauTN[id=C4-B12-D01-TH-044,muc=TH,loigiai={$\displaystyle\int_{2}^{6}f'(x)\,\mathrm{d}x=f(6)-f(2)=13$.}]{Cho hàm số $f(x)$ có đạo hàm liên tục trên $\mathbb{R}$, $f(2)=-5$ và $f(6)=8$. Giá trị của $\displaystyle\int_{2}^{6}f'(x)\,\mathrm{d}x$ bằng}
  {$13$}
  {$-13$}{$3$}{$8$}

\cauTN[id=C4-B12-D01-TH-045,muc=TH,loigiai={$\displaystyle\int_{0}^{2}\left[f(x)+2x+1\right]\,\mathrm{d}x=1\displaystyle\int_{0}^{2}f(x)\,\mathrm{d}x+\displaystyle\int_{0}^{2}(2x+1)\,\mathrm{d}x=1\cdot(-7)+6=-1$.}]{Biết $\displaystyle\int_{0}^{2}f(x)\,\mathrm{d}x=-7$. Giá trị của $\displaystyle\int_{0}^{2}\left[f(x)+2x+1\right]\,\mathrm{d}x$ bằng}
  {$-1$}
  {$3$}{$-4$}{$-2$}

\cauTN[id=C4-B12-D01-TH-046,muc=TH,loigiai={$\displaystyle\int_{-1}^{0}f'(x)\,\mathrm{d}x=f(0)-f(-1)=5$.}]{Cho hàm số $f(x)$ có đạo hàm liên tục trên $\mathbb{R}$, $f(-1)=0$ và $f(0)=5$. Giá trị của $\displaystyle\int_{-1}^{0}f'(x)\,\mathrm{d}x$ bằng}
  {$5$}
  {$-5$}{$0$}{$10$}

\cauTN[id=C4-B12-D01-TH-047,muc=TH,loigiai={$\displaystyle\int_{-1}^{3}f'(x)\,\mathrm{d}x=f(3)-f(-1)=3$.}]{Cho hàm số $f(x)$ có đạo hàm liên tục trên $\mathbb{R}$, $f(-1)=1$ và $f(3)=4$. Giá trị của $\displaystyle\int_{-1}^{3}f'(x)\,\mathrm{d}x$ bằng}
  {$3$}
  {$-3$}{$5$}{$4$}

\cauTN[id=C4-B12-D01-TH-048,muc=TH,loigiai={$\displaystyle\int_{5}^{6}f(x)\,\mathrm{d}x=\displaystyle\int_{1}^{6}f(x)\,\mathrm{d}x-\displaystyle\int_{1}^{5}f(x)\,\mathrm{d}x=-6$ nên $\displaystyle\int_{5}^{6}-3 f(x)\,\mathrm{d}x=18$.}]{Cho hàm số $f(x)$ liên tục trên $\mathbb{R}$, $\displaystyle\int_{1}^{5}f(x)\,\mathrm{d}x=4$ và $\displaystyle\int_{1}^{6}f(x)\,\mathrm{d}x=-2$. Giá trị của $\displaystyle\int_{5}^{6}-3 f(x)\,\mathrm{d}x$ bằng}
  {$18$}
  {$-6$}{$-18$}{$6$}

\cauDS[id=C4-B12-D01-TH-049,muc=TH,loigiai={Áp dụng tính chất tuyến tính và tính chất đổi cận của tích phân; $\displaystyle\int_{-2}^{1}x\,\mathrm{d}x=-\dfrac{3}{2}$. a) Sai: không có tính chất tích phân của tích. b) Sai. c) Sai: đúng là $-\dfrac{11}{2}$. d) Đúng.}]{Cho các hàm số $f(x)$, $g(x)$ liên tục trên $\mathbb{R}$ thoả mãn $\displaystyle\int_{-2}^{1}f(x)\,\mathrm{d}x=-4$ và $\displaystyle\int_{-2}^{1}g(x)\,\mathrm{d}x=-6$. Xét tính đúng sai của các mệnh đề sau:}{%
  \y{$\displaystyle\int_{-2}^{1}f(x)g(x)\,\mathrm{d}x=24$.}{s}
  \y{$\displaystyle\int_{-2}^{1}\left[-3f(x)-3g(x)\right]\,\mathrm{d}x=-6$.}{s}
  \y{$\displaystyle\int_{-2}^{1}\left[f(x)+x\right]\,\mathrm{d}x=-1$.}{s}
  \y{$\displaystyle\int_{-2}^{1}\left[f(x)+x\right]\,\mathrm{d}x=-\dfrac{11}{2}$.}{d}}

\cauDS[id=C4-B12-D01-TH-050,muc=TH,loigiai={Áp dụng tính chất tuyến tính và tính chất đổi cận của tích phân; $\displaystyle\int_{1}^{4}x\,\mathrm{d}x=\dfrac{15}{2}$. a) Đúng. b) Sai: đúng là $\dfrac{13}{2}$. c) Đúng. d) Đúng.}]{Cho các hàm số $f(x)$, $g(x)$ liên tục trên $\mathbb{R}$ thoả mãn $\displaystyle\int_{1}^{4}f(x)\,\mathrm{d}x=-1$ và $\displaystyle\int_{1}^{4}g(x)\,\mathrm{d}x=7$. Xét tính đúng sai của các mệnh đề sau:}{%
  \y{$\displaystyle\int_{1}^{4}\left[f(x)+x\right]\,\mathrm{d}x=\dfrac{13}{2}$.}{d}
  \y{$\displaystyle\int_{1}^{4}\left[f(x)+x\right]\,\mathrm{d}x=2$.}{s}
  \y{$\displaystyle\int_{1}^{4}\left[-2f(x)-2g(x)\right]\,\mathrm{d}x=-12$.}{d}
  \y{$\displaystyle\int_{1}^{4}\left[f(x)-g(x)\right]\,\mathrm{d}x=-8$.}{d}}

\cauDS[id=C4-B12-D01-TH-051,muc=TH,loigiai={Áp dụng tính chất tuyến tính và tính chất đổi cận của tích phân; $\displaystyle\int_{1}^{5}x\,\mathrm{d}x=12$. a) Sai. b) Sai: không có tính chất tích phân của tích. c) Sai: đúng là $8$. d) Đúng.}]{Cho các hàm số $f(x)$, $g(x)$ liên tục trên $\mathbb{R}$ thoả mãn $\displaystyle\int_{1}^{5}f(x)\,\mathrm{d}x=-4$ và $\displaystyle\int_{1}^{5}g(x)\,\mathrm{d}x=2$. Xét tính đúng sai của các mệnh đề sau:}{%
  \y{$\displaystyle\int_{1}^{5}\left[-2f(x)+3g(x)\right]\,\mathrm{d}x=2$.}{s}
  \y{$\displaystyle\int_{1}^{5}f(x)g(x)\,\mathrm{d}x=-8$.}{s}
  \y{$\displaystyle\int_{1}^{5}\left[f(x)+x\right]\,\mathrm{d}x=0$.}{s}
  \y{$\displaystyle\int_{1}^{5}\left[f(x)+x\right]\,\mathrm{d}x=8$.}{d}}

\cauDS[id=C4-B12-D01-TH-052,muc=TH,loigiai={Áp dụng tính chất tuyến tính và tính chất đổi cận của tích phân; $\displaystyle\int_{2}^{3}x\,\mathrm{d}x=\dfrac{5}{2}$. a) Sai. b) Đúng. c) Đúng. d) Sai: đúng là $\dfrac{15}{2}$.}]{Cho các hàm số $f(x)$, $g(x)$ liên tục trên $\mathbb{R}$ thoả mãn $\displaystyle\int_{2}^{3}f(x)\,\mathrm{d}x=5$ và $\displaystyle\int_{2}^{3}g(x)\,\mathrm{d}x=-7$. Xét tính đúng sai của các mệnh đề sau:}{%
  \y{$\displaystyle\int_{2}^{3}\left[2f(x)+2g(x)\right]\,\mathrm{d}x=24$.}{s}
  \y{$\displaystyle\int_{2}^{3}\left[2f(x)+2g(x)\right]\,\mathrm{d}x=-4$.}{d}
  \y{$\displaystyle\int_{2}^{3}\left[f(x)-g(x)\right]\,\mathrm{d}x=12$.}{d}
  \y{$\displaystyle\int_{2}^{3}\left[f(x)+x\right]\,\mathrm{d}x=6$.}{s}}

\cauDS[id=C4-B12-D01-TH-053,muc=TH,loigiai={Áp dụng tính chất tuyến tính và tính chất đổi cận của tích phân; $\displaystyle\int_{-2}^{-1}x\,\mathrm{d}x=-\dfrac{3}{2}$. a) Đúng. b) Sai: đúng là $\dfrac{3}{2}$. c) Sai. d) Đúng.}]{Cho các hàm số $f(x)$, $g(x)$ liên tục trên $\mathbb{R}$ thoả mãn $\displaystyle\int_{-2}^{-1}f(x)\,\mathrm{d}x=3$ và $\displaystyle\int_{-2}^{-1}g(x)\,\mathrm{d}x=3$. Xét tính đúng sai của các mệnh đề sau:}{%
  \y{$\displaystyle\int_{-2}^{-1}\left[f(x)+x\right]\,\mathrm{d}x=\dfrac{3}{2}$.}{d}
  \y{$\displaystyle\int_{-2}^{-1}\left[f(x)+x\right]\,\mathrm{d}x=4$.}{s}
  \y{$\displaystyle\int_{-2}^{-1}\left[3f(x)+2g(x)\right]\,\mathrm{d}x=3$.}{s}
  \y{$\displaystyle\int_{-1}^{-2}g(x)\,\mathrm{d}x=-3$.}{d}}

\cauDS[id=C4-B12-D01-TH-054,muc=TH,loigiai={Áp dụng tính chất tuyến tính và tính chất đổi cận của tích phân; $\displaystyle\int_{1}^{3}x\,\mathrm{d}x=4$. a) Sai. b) Sai: đúng là $13$. c) Đúng. d) Sai: không có tính chất tích phân của tích.}]{Cho các hàm số $f(x)$, $g(x)$ liên tục trên $\mathbb{R}$ thoả mãn $\displaystyle\int_{1}^{3}f(x)\,\mathrm{d}x=9$ và $\displaystyle\int_{1}^{3}g(x)\,\mathrm{d}x=-7$. Xét tính đúng sai của các mệnh đề sau:}{%
  \y{$\displaystyle\int_{1}^{3}\left[-f(x)+3g(x)\right]\,\mathrm{d}x=12$.}{s}
  \y{$\displaystyle\int_{1}^{3}\left[f(x)+x\right]\,\mathrm{d}x=11$.}{s}
  \y{$\displaystyle\int_{1}^{3}\left[-f(x)+3g(x)\right]\,\mathrm{d}x=-30$.}{d}
  \y{$\displaystyle\int_{1}^{3}f(x)g(x)\,\mathrm{d}x=-63$.}{s}}

\cauDS[id=C4-B12-D01-TH-055,muc=TH,loigiai={Áp dụng tính chất tuyến tính và tính chất đổi cận của tích phân; $\displaystyle\int_{1}^{4}x\,\mathrm{d}x=\dfrac{15}{2}$. a) Sai: đúng là $\dfrac{9}{2}$. b) Đúng. c) Đúng. d) Sai: không có tính chất tích phân của tích.}]{Cho các hàm số $f(x)$, $g(x)$ liên tục trên $\mathbb{R}$ thoả mãn $\displaystyle\int_{1}^{4}f(x)\,\mathrm{d}x=-3$ và $\displaystyle\int_{1}^{4}g(x)\,\mathrm{d}x=-4$. Xét tính đúng sai của các mệnh đề sau:}{%
  \y{$\displaystyle\int_{1}^{4}\left[f(x)+x\right]\,\mathrm{d}x=0$.}{s}
  \y{$\displaystyle\int_{1}^{4}\left[f(x)-g(x)\right]\,\mathrm{d}x=1$.}{d}
  \y{$\displaystyle\int_{1}^{4}\left[f(x)-2g(x)\right]\,\mathrm{d}x=5$.}{d}
  \y{$\displaystyle\int_{1}^{4}f(x)g(x)\,\mathrm{d}x=12$.}{s}}

\cauDS[id=C4-B12-D01-TH-056,muc=TH,loigiai={Áp dụng tính chất tuyến tính và tính chất đổi cận của tích phân; $\displaystyle\int_{2}^{4}x\,\mathrm{d}x=6$. a) Đúng. b) Đúng. c) Đúng. d) Sai: không có tính chất tích phân của tích.}]{Cho các hàm số $f(x)$, $g(x)$ liên tục trên $\mathbb{R}$ thoả mãn $\displaystyle\int_{2}^{4}f(x)\,\mathrm{d}x=8$ và $\displaystyle\int_{2}^{4}g(x)\,\mathrm{d}x=-6$. Xét tính đúng sai của các mệnh đề sau:}{%
  \y{$\displaystyle\int_{4}^{2}g(x)\,\mathrm{d}x=6$.}{d}
  \y{$\displaystyle\int_{2}^{4}\left[f(x)+g(x)\right]\,\mathrm{d}x=2$.}{d}
  \y{$\displaystyle\int_{2}^{4}\left[f(x)-g(x)\right]\,\mathrm{d}x=14$.}{d}
  \y{$\displaystyle\int_{2}^{4}f(x)g(x)\,\mathrm{d}x=-48$.}{s}}

\cauDS[id=C4-B12-D01-TH-057,muc=TH,loigiai={Áp dụng tính chất tuyến tính và tính chất đổi cận của tích phân; $\displaystyle\int_{2}^{4}x\,\mathrm{d}x=6$. a) Đúng. b) Đúng. c) Sai: đúng là $10$. d) Sai.}]{Cho các hàm số $f(x)$, $g(x)$ liên tục trên $\mathbb{R}$ thoả mãn $\displaystyle\int_{2}^{4}f(x)\,\mathrm{d}x=4$ và $\displaystyle\int_{2}^{4}g(x)\,\mathrm{d}x=7$. Xét tính đúng sai của các mệnh đề sau:}{%
  \y{$\displaystyle\int_{2}^{4}\left[-f(x)+2g(x)\right]\,\mathrm{d}x=10$.}{d}
  \y{$\displaystyle\int_{2}^{4}\left[f(x)-g(x)\right]\,\mathrm{d}x=-3$.}{d}
  \y{$\displaystyle\int_{2}^{4}\left[f(x)+x\right]\,\mathrm{d}x=6$.}{s}
  \y{$\displaystyle\int_{2}^{4}\left[-f(x)+2g(x)\right]\,\mathrm{d}x=-18$.}{s}}

\cauDS[id=C4-B12-D01-TH-058,muc=TH,loigiai={Áp dụng tính chất tuyến tính và tính chất đổi cận của tích phân; $\displaystyle\int_{2}^{3}x\,\mathrm{d}x=\dfrac{5}{2}$. a) Đúng. b) Sai: đúng là $-\dfrac{7}{2}$. c) Sai: không có tính chất tích phân của tích. d) Sai.}]{Cho các hàm số $f(x)$, $g(x)$ liên tục trên $\mathbb{R}$ thoả mãn $\displaystyle\int_{2}^{3}f(x)\,\mathrm{d}x=-6$ và $\displaystyle\int_{2}^{3}g(x)\,\mathrm{d}x=-5$. Xét tính đúng sai của các mệnh đề sau:}{%
  \y{$\displaystyle\int_{2}^{3}\left[f(x)+3g(x)\right]\,\mathrm{d}x=-21$.}{d}
  \y{$\displaystyle\int_{2}^{3}\left[f(x)+x\right]\,\mathrm{d}x=-5$.}{s}
  \y{$\displaystyle\int_{2}^{3}f(x)g(x)\,\mathrm{d}x=30$.}{s}
  \y{$\displaystyle\int_{2}^{3}\left[f(x)+3g(x)\right]\,\mathrm{d}x=9$.}{s}}

\cauDS[id=C4-B12-D01-TH-059,muc=TH,loigiai={Áp dụng tính chất tuyến tính và tính chất đổi cận của tích phân; $\displaystyle\int_{2}^{5}x\,\mathrm{d}x=\dfrac{21}{2}$. a) Sai: đúng là $\dfrac{15}{2}$. b) Sai. c) Đúng. d) Đúng.}]{Cho các hàm số $f(x)$, $g(x)$ liên tục trên $\mathbb{R}$ thoả mãn $\displaystyle\int_{2}^{5}f(x)\,\mathrm{d}x=-3$ và $\displaystyle\int_{2}^{5}g(x)\,\mathrm{d}x=6$. Xét tính đúng sai của các mệnh đề sau:}{%
  \y{$\displaystyle\int_{2}^{5}\left[f(x)+x\right]\,\mathrm{d}x=0$.}{s}
  \y{$\displaystyle\int_{2}^{5}\left[-3f(x)+g(x)\right]\,\mathrm{d}x=3$.}{s}
  \y{$\displaystyle\int_{2}^{5}\left[-3f(x)+g(x)\right]\,\mathrm{d}x=15$.}{d}
  \y{$\displaystyle\int_{5}^{2}g(x)\,\mathrm{d}x=-6$.}{d}}

\cauDS[id=C4-B12-D01-TH-060,muc=TH,loigiai={Áp dụng tính chất tuyến tính và tính chất đổi cận của tích phân; $\displaystyle\int_{0}^{3}x\,\mathrm{d}x=\dfrac{9}{2}$. a) Đúng. b) Đúng. c) Sai: đúng là $\dfrac{17}{2}$. d) Sai.}]{Cho các hàm số $f(x)$, $g(x)$ liên tục trên $\mathbb{R}$ thoả mãn $\displaystyle\int_{0}^{3}f(x)\,\mathrm{d}x=4$ và $\displaystyle\int_{0}^{3}g(x)\,\mathrm{d}x=7$. Xét tính đúng sai của các mệnh đề sau:}{%
  \y{$\displaystyle\int_{0}^{3}\left[f(x)-g(x)\right]\,\mathrm{d}x=-3$.}{d}
  \y{$\displaystyle\int_{0}^{3}\left[f(x)+x\right]\,\mathrm{d}x=\dfrac{17}{2}$.}{d}
  \y{$\displaystyle\int_{0}^{3}\left[f(x)+x\right]\,\mathrm{d}x=7$.}{s}
  \y{$\displaystyle\int_{0}^{3}\left[-3f(x)-g(x)\right]\,\mathrm{d}x=-5$.}{s}}

\cauDS[id=C4-B12-D01-TH-061,muc=TH,loigiai={Áp dụng tính chất tuyến tính và tính chất đổi cận của tích phân; $\displaystyle\int_{2}^{4}x\,\mathrm{d}x=6$. a) Sai: đúng là $12$. b) Đúng. c) Sai: không có tính chất tích phân của tích. d) Đúng.}]{Cho các hàm số $f(x)$, $g(x)$ liên tục trên $\mathbb{R}$ thoả mãn $\displaystyle\int_{2}^{4}f(x)\,\mathrm{d}x=6$ và $\displaystyle\int_{2}^{4}g(x)\,\mathrm{d}x=4$. Xét tính đúng sai của các mệnh đề sau:}{%
  \y{$\displaystyle\int_{2}^{4}\left[f(x)+x\right]\,\mathrm{d}x=8$.}{s}
  \y{$\displaystyle\int_{2}^{4}\left[f(x)+x\right]\,\mathrm{d}x=12$.}{d}
  \y{$\displaystyle\int_{2}^{4}f(x)g(x)\,\mathrm{d}x=24$.}{s}
  \y{$\displaystyle\int_{2}^{4}\left[f(x)+2g(x)\right]\,\mathrm{d}x=14$.}{d}}

\cauDS[id=C4-B12-D01-TH-062,muc=TH,loigiai={Áp dụng tính chất tuyến tính và tính chất đổi cận của tích phân; $\displaystyle\int_{0}^{4}x\,\mathrm{d}x=8$. a) Sai: đúng là $13$. b) Sai: không có tính chất tích phân của tích. c) Đúng. d) Đúng.}]{Cho các hàm số $f(x)$, $g(x)$ liên tục trên $\mathbb{R}$ thoả mãn $\displaystyle\int_{0}^{4}f(x)\,\mathrm{d}x=5$ và $\displaystyle\int_{0}^{4}g(x)\,\mathrm{d}x=7$. Xét tính đúng sai của các mệnh đề sau:}{%
  \y{$\displaystyle\int_{0}^{4}\left[f(x)+x\right]\,\mathrm{d}x=9$.}{s}
  \y{$\displaystyle\int_{0}^{4}f(x)g(x)\,\mathrm{d}x=35$.}{s}
  \y{$\displaystyle\int_{0}^{4}\left[-f(x)+3g(x)\right]\,\mathrm{d}x=16$.}{d}
  \y{$\displaystyle\int_{0}^{4}\left[f(x)+x\right]\,\mathrm{d}x=13$.}{d}}

\cauDS[id=C4-B12-D01-TH-063,muc=TH,loigiai={Áp dụng tính chất tuyến tính và tính chất đổi cận của tích phân; $\displaystyle\int_{-2}^{1}x\,\mathrm{d}x=-\dfrac{3}{2}$. a) Sai: đúng là $-\dfrac{5}{2}$. b) Đúng. c) Sai. d) Sai: không có tính chất tích phân của tích.}]{Cho các hàm số $f(x)$, $g(x)$ liên tục trên $\mathbb{R}$ thoả mãn $\displaystyle\int_{-2}^{1}f(x)\,\mathrm{d}x=-1$ và $\displaystyle\int_{-2}^{1}g(x)\,\mathrm{d}x=3$. Xét tính đúng sai của các mệnh đề sau:}{%
  \y{$\displaystyle\int_{-2}^{1}\left[f(x)+x\right]\,\mathrm{d}x=2$.}{s}
  \y{$\displaystyle\int_{-2}^{1}\left[f(x)-g(x)\right]\,\mathrm{d}x=-4$.}{d}
  \y{$\displaystyle\int_{-2}^{1}\left[f(x)-2g(x)\right]\,\mathrm{d}x=5$.}{s}
  \y{$\displaystyle\int_{-2}^{1}f(x)g(x)\,\mathrm{d}x=-3$.}{s}}

\cauDS[id=C4-B12-D01-TH-064,muc=TH,loigiai={Áp dụng tính chất tuyến tính và tính chất đổi cận của tích phân; $\displaystyle\int_{1}^{3}x\,\mathrm{d}x=4$. a) Đúng. b) Sai: không có tính chất tích phân của tích. c) Đúng. d) Sai: đúng là $3$.}]{Cho các hàm số $f(x)$, $g(x)$ liên tục trên $\mathbb{R}$ thoả mãn $\displaystyle\int_{1}^{3}f(x)\,\mathrm{d}x=-1$ và $\displaystyle\int_{1}^{3}g(x)\,\mathrm{d}x=9$. Xét tính đúng sai của các mệnh đề sau:}{%
  \y{$\displaystyle\int_{1}^{3}\left[f(x)+x\right]\,\mathrm{d}x=3$.}{d}
  \y{$\displaystyle\int_{1}^{3}f(x)g(x)\,\mathrm{d}x=-9$.}{s}
  \y{$\displaystyle\int_{3}^{1}g(x)\,\mathrm{d}x=-9$.}{d}
  \y{$\displaystyle\int_{1}^{3}\left[f(x)+x\right]\,\mathrm{d}x=1$.}{s}}

\cauTLN[id=C4-B12-D01-TH-065,muc=TH,loigiai={$\displaystyle\int_{0}^{2}\left[f(x)-x-3\right]\,\mathrm{d}x=1\cdot(-5)-8=-13$.}]{Biết $\displaystyle\int_{0}^{2}f(x)\,\mathrm{d}x=-5$. Tính $\displaystyle\int_{0}^{2}\left[f(x)-x-3\right]\,\mathrm{d}x$.}{-13}

\cauTLN[id=C4-B12-D01-TH-066,muc=TH,loigiai={$\displaystyle\int_{2}^{4}\left[-2f(x)-2x-1\right]\,\mathrm{d}x=-2\cdot(-3)-14=-8$.}]{Biết $\displaystyle\int_{2}^{4}f(x)\,\mathrm{d}x=-3$. Tính $\displaystyle\int_{2}^{4}\left[-2f(x)-2x-1\right]\,\mathrm{d}x$.}{-8}

\cauTLN[id=C4-B12-D01-TH-067,muc=TH,loigiai={$\displaystyle\int_{1}^{5}\left[2f(x)-x+3\right]\,\mathrm{d}x=2\cdot(4)+0=8$.}]{Biết $\displaystyle\int_{1}^{5}f(x)\,\mathrm{d}x=4$. Tính $\displaystyle\int_{1}^{5}\left[2f(x)-x+3\right]\,\mathrm{d}x$.}{8}

\cauTLN[id=C4-B12-D01-TH-068,muc=TH,loigiai={$\displaystyle\int_{-1}^{0}\left[f(x)+2x+1\right]\,\mathrm{d}x=1\cdot(-7)+0=-7$.}]{Biết $\displaystyle\int_{-1}^{0}f(x)\,\mathrm{d}x=-7$. Tính $\displaystyle\int_{-1}^{0}\left[f(x)+2x+1\right]\,\mathrm{d}x$.}{-7}

\cauTLN[id=C4-B12-D01-TH-069,muc=TH,loigiai={$\displaystyle\int_{1}^{5}\left[4f(x)-x-2\right]\,\mathrm{d}x=4\cdot(9)-20=16$.}]{Biết $\displaystyle\int_{1}^{5}f(x)\,\mathrm{d}x=9$. Tính $\displaystyle\int_{1}^{5}\left[4f(x)-x-2\right]\,\mathrm{d}x$.}{16}

\cauTLN[id=C4-B12-D01-TH-070,muc=TH,loigiai={$\displaystyle\int_{-1}^{1}\left[-2f(x)+x+3\right]\,\mathrm{d}x=-2\cdot(-1)+6=8$.}]{Biết $\displaystyle\int_{-1}^{1}f(x)\,\mathrm{d}x=-1$. Tính $\displaystyle\int_{-1}^{1}\left[-2f(x)+x+3\right]\,\mathrm{d}x$.}{8}

\cauTLN[id=C4-B12-D01-TH-071,muc=TH,loigiai={$\displaystyle\int_{0}^{1}\left[-f(x)+x-2\right]\,\mathrm{d}x=-1\cdot(1)-\dfrac{3}{2}=-\dfrac{5}{2}$.}]{Biết $\displaystyle\int_{0}^{1}f(x)\,\mathrm{d}x=1$. Tính $\displaystyle\int_{0}^{1}\left[-f(x)+x-2\right]\,\mathrm{d}x$.}{-2,5}

\cauTLN[id=C4-B12-D01-TH-072,muc=TH,loigiai={$\displaystyle\int_{0}^{3}\left[3f(x)+x-4\right]\,\mathrm{d}x=3\cdot(7)-\dfrac{15}{2}=\dfrac{27}{2}$.}]{Biết $\displaystyle\int_{0}^{3}f(x)\,\mathrm{d}x=7$. Tính $\displaystyle\int_{0}^{3}\left[3f(x)+x-4\right]\,\mathrm{d}x$.}{13,5}

\cauTLN[id=C4-B12-D01-TH-073,muc=TH,loigiai={$\displaystyle\int_{1}^{3}\left[-2f(x)+2x-1\right]\,\mathrm{d}x=-2\cdot(-2)+6=10$.}]{Biết $\displaystyle\int_{1}^{3}f(x)\,\mathrm{d}x=-2$. Tính $\displaystyle\int_{1}^{3}\left[-2f(x)+2x-1\right]\,\mathrm{d}x$.}{10}

\cauTLN[id=C4-B12-D01-TH-074,muc=TH,loigiai={$\displaystyle\int_{2}^{3}\left[3f(x)+x-2\right]\,\mathrm{d}x=3\cdot(2)+\dfrac{1}{2}=\dfrac{13}{2}$.}]{Biết $\displaystyle\int_{2}^{3}f(x)\,\mathrm{d}x=2$. Tính $\displaystyle\int_{2}^{3}\left[3f(x)+x-2\right]\,\mathrm{d}x$.}{6,5}

\cauTLN[id=C4-B12-D01-TH-075,muc=TH,loigiai={$\displaystyle\int_{0}^{1}\left[f(x)+3x-1\right]\,\mathrm{d}x=1\cdot(-6)+\dfrac{1}{2}=-\dfrac{11}{2}$.}]{Biết $\displaystyle\int_{0}^{1}f(x)\,\mathrm{d}x=-6$. Tính $\displaystyle\int_{0}^{1}\left[f(x)+3x-1\right]\,\mathrm{d}x$.}{-5,5}

\cauTLN[id=C4-B12-D01-TH-076,muc=TH,loigiai={$\displaystyle\int_{-1}^{1}\left[4f(x)-x-2\right]\,\mathrm{d}x=4\cdot(-4)-4=-20$.}]{Biết $\displaystyle\int_{-1}^{1}f(x)\,\mathrm{d}x=-4$. Tính $\displaystyle\int_{-1}^{1}\left[4f(x)-x-2\right]\,\mathrm{d}x$.}{-20}

\cauTLN[id=C4-B12-D01-TH-077,muc=TH,loigiai={$\displaystyle\int_{1}^{5}\left[-2f(x)-2x-2\right]\,\mathrm{d}x=-2\cdot(-2)-32=-28$.}]{Biết $\displaystyle\int_{1}^{5}f(x)\,\mathrm{d}x=-2$. Tính $\displaystyle\int_{1}^{5}\left[-2f(x)-2x-2\right]\,\mathrm{d}x$.}{-28}

\cauTLN[id=C4-B12-D01-TH-078,muc=TH,loigiai={$\displaystyle\int_{0}^{4}\left[f(x)+x-1\right]\,\mathrm{d}x=1\cdot(4)+4=8$.}]{Biết $\displaystyle\int_{0}^{4}f(x)\,\mathrm{d}x=4$. Tính $\displaystyle\int_{0}^{4}\left[f(x)+x-1\right]\,\mathrm{d}x$.}{8}

\cauTLN[id=C4-B12-D01-TH-079,muc=TH,loigiai={$\displaystyle\int_{0}^{3}\left[4f(x)+2x-1\right]\,\mathrm{d}x=4\cdot(-4)+6=-10$.}]{Biết $\displaystyle\int_{0}^{3}f(x)\,\mathrm{d}x=-4$. Tính $\displaystyle\int_{0}^{3}\left[4f(x)+2x-1\right]\,\mathrm{d}x$.}{-10}

\cauTLN[id=C4-B12-D01-TH-080,muc=TH,loigiai={$\displaystyle\int_{1}^{2}\left[f(x)+3x+3\right]\,\mathrm{d}x=1\cdot(9)+\dfrac{15}{2}=\dfrac{33}{2}$.}]{Biết $\displaystyle\int_{1}^{2}f(x)\,\mathrm{d}x=9$. Tính $\displaystyle\int_{1}^{2}\left[f(x)+3x+3\right]\,\mathrm{d}x$.}{16,5}

\cauTN[id=C4-B12-D01-VD-001,muc=VD,loigiai={$\displaystyle\int_{0}^{1}\left[2f'(x)-2x\right]\,\mathrm{d}x=2\left[f(1)-f(0)\right]+\displaystyle\int_{0}^{1}-2x\,\mathrm{d}x=2\cdot-6-1=-13$.}]{Cho hàm số $f(x)$ có đạo hàm liên tục trên $\mathbb{R}$, $f(0)=1$ và $f(1)=-5$. Giá trị của $\displaystyle\int_{0}^{1}\left[2f'(x)-2x\right]\,\mathrm{d}x$ bằng}
  {$-13$}
  {$-7$}{$-14$}{$-11$}

\cauTN[id=C4-B12-D01-VD-002,muc=VD,loigiai={$\displaystyle\int_{1}^{2}\left[2f'(x)+3x\right]\,\mathrm{d}x=2\left[f(2)-f(1)\right]+\displaystyle\int_{1}^{2}3x\,\mathrm{d}x=2\cdot2+\dfrac{9}{2}=\dfrac{17}{2}$.}]{Cho hàm số $f(x)$ có đạo hàm liên tục trên $\mathbb{R}$, $f(1)=-5$ và $f(2)=-3$. Giá trị của $\displaystyle\int_{1}^{2}\left[2f'(x)+3x\right]\,\mathrm{d}x$ bằng}
  {$\dfrac{17}{2}$}
  {$\dfrac{13}{2}$}{$7$}{$-\dfrac{3}{2}$}

\cauTN[id=C4-B12-D01-VD-003,muc=VD,loigiai={$\displaystyle\int_{0}^{1}\left[2f'(x)+x\right]\,\mathrm{d}x=2\left[f(1)-f(0)\right]+\displaystyle\int_{0}^{1}x\,\mathrm{d}x=2\cdot5+\dfrac{1}{2}=\dfrac{21}{2}$.}]{Cho hàm số $f(x)$ có đạo hàm liên tục trên $\mathbb{R}$, $f(0)=1$ và $f(1)=6$. Giá trị của $\displaystyle\int_{0}^{1}\left[2f'(x)+x\right]\,\mathrm{d}x$ bằng}
  {$\dfrac{21}{2}$}
  {$\dfrac{11}{2}$}{$11$}{$\dfrac{25}{2}$}

\cauTN[id=C4-B12-D01-VD-004,muc=VD,loigiai={$\displaystyle\int_{2}^{4}f(x)\,\mathrm{d}x=6-(1)=5$; $\displaystyle\int_{2}^{4}6x^2\,\mathrm{d}x=112$. Vậy $\displaystyle\int_{2}^{4}\left[2f(x)-6x^{2}\right]\,\mathrm{d}x=2\cdot(5)-(112)=-102$.}]{Cho hàm số $f(x)$ liên tục trên $\mathbb{R}$, $\displaystyle\int_{0}^{2}f(x)\,\mathrm{d}x=1$ và $\displaystyle\int_{0}^{4}f(x)\,\mathrm{d}x=6$. Giá trị của $\displaystyle\int_{2}^{4}\left[2f(x)-6x^{2}\right]\,\mathrm{d}x$ bằng}
  {$-102$}
  {$-98$}{$122$}{$-107$}

\cauTN[id=C4-B12-D01-VD-005,muc=VD,loigiai={Đặt $A=\displaystyle\int_{0}^{1}f(x)\,\mathrm{d}x$, $B=\displaystyle\int_{0}^{1}g(x)\,\mathrm{d}x$. Ta có hệ $A+B=-8$, $3A+B=-18$, suy ra $A=-5$, $B=-3$. Vậy giá trị cần tìm là $-2$.}]{Cho các hàm số $f(x)$, $g(x)$ liên tục trên $\mathbb{R}$ thoả mãn $\displaystyle\int_{0}^{1}\left[f(x)+g(x)\right]\,\mathrm{d}x=-8$ và $\displaystyle\int_{0}^{1}\left[3f(x)+g(x)\right]\,\mathrm{d}x=-18$. Giá trị của $\displaystyle\int_{0}^{1}\left[f(x)-g(x)\right]\,\mathrm{d}x$ bằng}
  {$-2$}
  {$-5$}{$2$}{$10$}

\cauTN[id=C4-B12-D01-VD-006,muc=VD,loigiai={$\displaystyle\int_{-1}^{0}\left[2f'(x)-2x\right]\,\mathrm{d}x=2\left[f(0)-f(-1)\right]+\displaystyle\int_{-1}^{0}-2x\,\mathrm{d}x=2\cdot3+1=7$.}]{Cho hàm số $f(x)$ có đạo hàm liên tục trên $\mathbb{R}$, $f(-1)=-5$ và $f(0)=-2$. Giá trị của $\displaystyle\int_{-1}^{0}\left[2f'(x)-2x\right]\,\mathrm{d}x$ bằng}
  {$7$}
  {$4$}{$-3$}{$5$}

\cauTN[id=C4-B12-D01-VD-007,muc=VD,loigiai={Đặt $A=\displaystyle\int_{0}^{2}f(x)\,\mathrm{d}x$, $B=\displaystyle\int_{0}^{2}g(x)\,\mathrm{d}x$. Ta có hệ $A+B=5$, $3A-2B=-15$, suy ra $A=-1$, $B=6$. Vậy giá trị cần tìm là $-7$.}]{Cho các hàm số $f(x)$, $g(x)$ liên tục trên $\mathbb{R}$ thoả mãn $\displaystyle\int_{0}^{2}\left[f(x)+g(x)\right]\,\mathrm{d}x=5$ và $\displaystyle\int_{0}^{2}\left[3f(x)-2g(x)\right]\,\mathrm{d}x=-15$. Giá trị của $\displaystyle\int_{0}^{2}\left[f(x)-g(x)\right]\,\mathrm{d}x$ bằng}
  {$-7$}
  {$-1$}{$7$}{$20$}

\cauTN[id=C4-B12-D01-VD-008,muc=VD,loigiai={Đặt $A=\displaystyle\int_{0}^{4}f(x)\,\mathrm{d}x$, $B=\displaystyle\int_{0}^{4}g(x)\,\mathrm{d}x$. Ta có hệ $A+B=1$, $3A-B=-13$, suy ra $A=-3$, $B=4$. Vậy giá trị cần tìm là $-3$.}]{Cho các hàm số $f(x)$, $g(x)$ liên tục trên $\mathbb{R}$ thoả mãn $\displaystyle\int_{0}^{4}\left[f(x)+g(x)\right]\,\mathrm{d}x=1$ và $\displaystyle\int_{0}^{4}\left[3f(x)-g(x)\right]\,\mathrm{d}x=-13$. Giá trị của $\displaystyle\int_{0}^{4}f(x)\,\mathrm{d}x$ bằng}
  {$-3$}
  {$4$}{$3$}{$14$}

\cauTN[id=C4-B12-D01-VD-009,muc=VD,loigiai={$\displaystyle\int_{2}^{3}f(x)\,\mathrm{d}x=8-(-3)=11$; $\displaystyle\int_{2}^{3}-6x^2\,\mathrm{d}x=-38$. Vậy $\displaystyle\int_{2}^{3}\left[2f(x)+6x^{2}\right]\,\mathrm{d}x=2\cdot(11)-(-38)=60$.}]{Cho hàm số $f(x)$ liên tục trên $\mathbb{R}$, $\displaystyle\int_{0}^{2}f(x)\,\mathrm{d}x=-3$ và $\displaystyle\int_{0}^{3}f(x)\,\mathrm{d}x=8$. Giá trị của $\displaystyle\int_{2}^{3}\left[2f(x)+6x^{2}\right]\,\mathrm{d}x$ bằng}
  {$60$}
  {$48$}{$-16$}{$49$}

\cauTN[id=C4-B12-D01-VD-010,muc=VD,loigiai={Đặt $A=\displaystyle\int_{0}^{4}f(x)\,\mathrm{d}x$, $B=\displaystyle\int_{0}^{4}g(x)\,\mathrm{d}x$. Ta có hệ $A+B=3$, $3A-2B=-11$, suy ra $A=-1$, $B=4$. Vậy giá trị cần tìm là $-1$.}]{Cho các hàm số $f(x)$, $g(x)$ liên tục trên $\mathbb{R}$ thoả mãn $\displaystyle\int_{0}^{4}\left[f(x)+g(x)\right]\,\mathrm{d}x=3$ và $\displaystyle\int_{0}^{4}\left[3f(x)-2g(x)\right]\,\mathrm{d}x=-11$. Giá trị của $\displaystyle\int_{0}^{4}f(x)\,\mathrm{d}x$ bằng}
  {$-1$}
  {$4$}{$1$}{$14$}

\cauTN[id=C4-B12-D01-VD-011,muc=VD,loigiai={$\displaystyle\int_{3}^{4}f(x)\,\mathrm{d}x=-5-(-1)=-4$; $\displaystyle\int_{3}^{4}-3x^2\,\mathrm{d}x=-37$. Vậy $\displaystyle\int_{3}^{4}\left[2f(x)+3x^{2}\right]\,\mathrm{d}x=2\cdot(-4)-(-37)=29$.}]{Cho hàm số $f(x)$ liên tục trên $\mathbb{R}$, $\displaystyle\int_{0}^{3}f(x)\,\mathrm{d}x=-1$ và $\displaystyle\int_{0}^{4}f(x)\,\mathrm{d}x=-5$. Giá trị của $\displaystyle\int_{3}^{4}\left[2f(x)+3x^{2}\right]\,\mathrm{d}x$ bằng}
  {$29$}
  {$25$}{$-45$}{$33$}

\cauTN[id=C4-B12-D01-VD-012,muc=VD,loigiai={$\displaystyle\int_{-1}^{2}\left[2f'(x)-3x\right]\,\mathrm{d}x=2\left[f(2)-f(-1)\right]+\displaystyle\int_{-1}^{2}-3x\,\mathrm{d}x=2\cdot-3-\dfrac{9}{2}=-\dfrac{21}{2}$.}]{Cho hàm số $f(x)$ có đạo hàm liên tục trên $\mathbb{R}$, $f(-1)=3$ và $f(2)=0$. Giá trị của $\displaystyle\int_{-1}^{2}\left[2f'(x)-3x\right]\,\mathrm{d}x$ bằng}
  {$-\dfrac{21}{2}$}
  {$-\dfrac{15}{2}$}{$-15$}{$-\dfrac{9}{2}$}

\cauTN[id=C4-B12-D01-VD-013,muc=VD,loigiai={$\displaystyle\int_{1}^{3}\left[2f'(x)-2x\right]\,\mathrm{d}x=2\left[f(3)-f(1)\right]+\displaystyle\int_{1}^{3}-2x\,\mathrm{d}x=2\cdot1-8=-6$.}]{Cho hàm số $f(x)$ có đạo hàm liên tục trên $\mathbb{R}$, $f(1)=-1$ và $f(3)=0$. Giá trị của $\displaystyle\int_{1}^{3}\left[2f'(x)-2x\right]\,\mathrm{d}x$ bằng}
  {$-6$}
  {$-7$}{$-2$}{$-8$}

\cauTN[id=C4-B12-D01-VD-014,muc=VD,loigiai={$\displaystyle\int_{-1}^{2}\left[2f'(x)+2x\right]\,\mathrm{d}x=2\left[f(2)-f(-1)\right]+\displaystyle\int_{-1}^{2}2x\,\mathrm{d}x=2\cdot5+3=13$.}]{Cho hàm số $f(x)$ có đạo hàm liên tục trên $\mathbb{R}$, $f(-1)=-2$ và $f(2)=3$. Giá trị của $\displaystyle\int_{-1}^{2}\left[2f'(x)+2x\right]\,\mathrm{d}x$ bằng}
  {$13$}
  {$8$}{$16$}{$9$}

\cauTN[id=C4-B12-D01-VD-015,muc=VD,loigiai={$\displaystyle\int_{2}^{3}f(x)\,\mathrm{d}x=-5-(6)=-11$; $\displaystyle\int_{2}^{3}6x^2\,\mathrm{d}x=38$. Vậy $\displaystyle\int_{2}^{3}\left[2f(x)-6x^{2}\right]\,\mathrm{d}x=2\cdot(-11)-(38)=-60$.}]{Cho hàm số $f(x)$ liên tục trên $\mathbb{R}$, $\displaystyle\int_{0}^{2}f(x)\,\mathrm{d}x=6$ và $\displaystyle\int_{0}^{3}f(x)\,\mathrm{d}x=-5$. Giá trị của $\displaystyle\int_{2}^{3}\left[2f(x)-6x^{2}\right]\,\mathrm{d}x$ bằng}
  {$-60$}
  {$-36$}{$16$}{$-49$}

\cauTN[id=C4-B12-D01-VD-016,muc=VD,loigiai={$\displaystyle\int_{1}^{2}f(x)\,\mathrm{d}x=3-(3)=0$; $\displaystyle\int_{1}^{2}-6x^2\,\mathrm{d}x=-14$. Vậy $\displaystyle\int_{1}^{2}\left[2f(x)+6x^{2}\right]\,\mathrm{d}x=2\cdot(0)-(-14)=14$.}]{Cho hàm số $f(x)$ liên tục trên $\mathbb{R}$, $\displaystyle\int_{0}^{1}f(x)\,\mathrm{d}x=3$ và $\displaystyle\int_{0}^{2}f(x)\,\mathrm{d}x=3$. Giá trị của $\displaystyle\int_{1}^{2}\left[2f(x)+6x^{2}\right]\,\mathrm{d}x$ bằng}
  {$14$}
  {$26$}{$-14$}{$6$}

\cauTN[id=C4-B12-D01-VD-017,muc=VD,loigiai={$\displaystyle\int_{-1}^{2}\left[2f'(x)-3x\right]\,\mathrm{d}x=2\left[f(2)-f(-1)\right]+\displaystyle\int_{-1}^{2}-3x\,\mathrm{d}x=2\cdot1-\dfrac{9}{2}=-\dfrac{5}{2}$.}]{Cho hàm số $f(x)$ có đạo hàm liên tục trên $\mathbb{R}$, $f(-1)=3$ và $f(2)=4$. Giá trị của $\displaystyle\int_{-1}^{2}\left[2f'(x)-3x\right]\,\mathrm{d}x$ bằng}
  {$-\dfrac{5}{2}$}
  {$-\dfrac{7}{2}$}{$-7$}{$\dfrac{7}{2}$}

\cauTN[id=C4-B12-D01-VD-018,muc=VD,loigiai={$\displaystyle\int_{0}^{1}\left[2f'(x)-x\right]\,\mathrm{d}x=2\left[f(1)-f(0)\right]+\displaystyle\int_{0}^{1}-x\,\mathrm{d}x=2\cdot-4-\dfrac{1}{2}=-\dfrac{17}{2}$.}]{Cho hàm số $f(x)$ có đạo hàm liên tục trên $\mathbb{R}$, $f(0)=3$ và $f(1)=-1$. Giá trị của $\displaystyle\int_{0}^{1}\left[2f'(x)-x\right]\,\mathrm{d}x$ bằng}
  {$-\dfrac{17}{2}$}
  {$-\dfrac{9}{2}$}{$-9$}{$-\dfrac{5}{2}$}

\cauTN[id=C4-B12-D01-VD-019,muc=VD,loigiai={Đặt $A=\displaystyle\int_{0}^{1}f(x)\,\mathrm{d}x$, $B=\displaystyle\int_{0}^{1}g(x)\,\mathrm{d}x$. Ta có hệ $A+B=0$, $2A+B=-4$, suy ra $A=-4$, $B=4$. Vậy giá trị cần tìm là $-8$.}]{Cho các hàm số $f(x)$, $g(x)$ liên tục trên $\mathbb{R}$ thoả mãn $\displaystyle\int_{0}^{1}\left[f(x)+g(x)\right]\,\mathrm{d}x=0$ và $\displaystyle\int_{0}^{1}\left[2f(x)+g(x)\right]\,\mathrm{d}x=-4$. Giá trị của $\displaystyle\int_{0}^{1}\left[f(x)-g(x)\right]\,\mathrm{d}x$ bằng}
  {$-8$}
  {$-4$}{$8$}{$4$}

\cauTN[id=C4-B12-D01-VD-020,muc=VD,loigiai={$\displaystyle\int_{2}^{3}f(x)\,\mathrm{d}x=-6-(6)=-12$; $\displaystyle\int_{2}^{3}3x^2\,\mathrm{d}x=19$. Vậy $\displaystyle\int_{2}^{3}\left[2f(x)-3x^{2}\right]\,\mathrm{d}x=2\cdot(-12)-(19)=-43$.}]{Cho hàm số $f(x)$ liên tục trên $\mathbb{R}$, $\displaystyle\int_{0}^{2}f(x)\,\mathrm{d}x=6$ và $\displaystyle\int_{0}^{3}f(x)\,\mathrm{d}x=-6$. Giá trị của $\displaystyle\int_{2}^{3}\left[2f(x)-3x^{2}\right]\,\mathrm{d}x$ bằng}
  {$-43$}
  {$-19$}{$-5$}{$-31$}

\cauTN[id=C4-B12-D01-VD-021,muc=VD,loigiai={Đặt $A=\displaystyle\int_{0}^{4}f(x)\,\mathrm{d}x$, $B=\displaystyle\int_{0}^{4}g(x)\,\mathrm{d}x$. Ta có hệ $A+B=-5$, $2A+B=-7$, suy ra $A=-2$, $B=-3$. Vậy giá trị cần tìm là $-3$.}]{Cho các hàm số $f(x)$, $g(x)$ liên tục trên $\mathbb{R}$ thoả mãn $\displaystyle\int_{0}^{4}\left[f(x)+g(x)\right]\,\mathrm{d}x=-5$ và $\displaystyle\int_{0}^{4}\left[2f(x)+g(x)\right]\,\mathrm{d}x=-7$. Giá trị của $\displaystyle\int_{0}^{4}g(x)\,\mathrm{d}x$ bằng}
  {$-3$}
  {$-2$}{$3$}{$2$}

\cauTN[id=C4-B12-D01-VD-022,muc=VD,loigiai={$\displaystyle\int_{1}^{3}\left[2f'(x)+3x\right]\,\mathrm{d}x=2\left[f(3)-f(1)\right]+\displaystyle\int_{1}^{3}3x\,\mathrm{d}x=2\cdot3+12=18$.}]{Cho hàm số $f(x)$ có đạo hàm liên tục trên $\mathbb{R}$, $f(1)=-5$ và $f(3)=-2$. Giá trị của $\displaystyle\int_{1}^{3}\left[2f'(x)+3x\right]\,\mathrm{d}x$ bằng}
  {$18$}
  {$15$}{$12$}{$8$}

\cauTN[id=C4-B12-D01-VD-023,muc=VD,loigiai={$\displaystyle\int_{1}^{2}f(x)\,\mathrm{d}x=-2-(-5)=3$; $\displaystyle\int_{1}^{2}-6x^2\,\mathrm{d}x=-14$. Vậy $\displaystyle\int_{1}^{2}\left[2f(x)+6x^{2}\right]\,\mathrm{d}x=2\cdot(3)-(-14)=20$.}]{Cho hàm số $f(x)$ liên tục trên $\mathbb{R}$, $\displaystyle\int_{0}^{1}f(x)\,\mathrm{d}x=-5$ và $\displaystyle\int_{0}^{2}f(x)\,\mathrm{d}x=-2$. Giá trị của $\displaystyle\int_{1}^{2}\left[2f(x)+6x^{2}\right]\,\mathrm{d}x$ bằng}
  {$20$}
  {$0$}{$-8$}{$17$}

\cauTN[id=C4-B12-D01-VD-024,muc=VD,loigiai={$\displaystyle\int_{0}^{3}\left[2f'(x)+2x\right]\,\mathrm{d}x=2\left[f(3)-f(0)\right]+\displaystyle\int_{0}^{3}2x\,\mathrm{d}x=2\cdot-1+9=7$.}]{Cho hàm số $f(x)$ có đạo hàm liên tục trên $\mathbb{R}$, $f(0)=-5$ và $f(3)=-6$. Giá trị của $\displaystyle\int_{0}^{3}\left[2f'(x)+2x\right]\,\mathrm{d}x$ bằng}
  {$7$}
  {$8$}{$4$}{$-3$}

\cauTN[id=C4-B12-D01-VD-025,muc=VD,loigiai={$\displaystyle\int_{0}^{1}\left[2f'(x)-3x\right]\,\mathrm{d}x=2\left[f(1)-f(0)\right]+\displaystyle\int_{0}^{1}-3x\,\mathrm{d}x=2\cdot-3-\dfrac{3}{2}=-\dfrac{15}{2}$.}]{Cho hàm số $f(x)$ có đạo hàm liên tục trên $\mathbb{R}$, $f(0)=3$ và $f(1)=0$. Giá trị của $\displaystyle\int_{0}^{1}\left[2f'(x)-3x\right]\,\mathrm{d}x$ bằng}
  {$-\dfrac{15}{2}$}
  {$-\dfrac{9}{2}$}{$-9$}{$-\dfrac{3}{2}$}

\cauTN[id=C4-B12-D01-VD-026,muc=VD,loigiai={$\displaystyle\int_{1}^{3}f(x)\,\mathrm{d}x=-1-(8)=-9$; $\displaystyle\int_{1}^{3}-3x^2\,\mathrm{d}x=-26$. Vậy $\displaystyle\int_{1}^{3}\left[2f(x)+3x^{2}\right]\,\mathrm{d}x=2\cdot(-9)-(-26)=8$.}]{Cho hàm số $f(x)$ liên tục trên $\mathbb{R}$, $\displaystyle\int_{0}^{1}f(x)\,\mathrm{d}x=8$ và $\displaystyle\int_{0}^{3}f(x)\,\mathrm{d}x=-1$. Giá trị của $\displaystyle\int_{1}^{3}\left[2f(x)+3x^{2}\right]\,\mathrm{d}x$ bằng}
  {$8$}
  {$40$}{$-44$}{$17$}

\cauTN[id=C4-B12-D01-VD-027,muc=VD,loigiai={$\displaystyle\int_{-1}^{0}\left[2f'(x)-3x\right]\,\mathrm{d}x=2\left[f(0)-f(-1)\right]+\displaystyle\int_{-1}^{0}-3x\,\mathrm{d}x=2\cdot2+\dfrac{3}{2}=\dfrac{11}{2}$.}]{Cho hàm số $f(x)$ có đạo hàm liên tục trên $\mathbb{R}$, $f(-1)=-4$ và $f(0)=-2$. Giá trị của $\displaystyle\int_{-1}^{0}\left[2f'(x)-3x\right]\,\mathrm{d}x$ bằng}
  {$\dfrac{11}{2}$}
  {$\dfrac{7}{2}$}{$1$}{$-\dfrac{5}{2}$}

\cauTN[id=C4-B12-D01-VD-028,muc=VD,loigiai={Đặt $A=\displaystyle\int_{0}^{1}f(x)\,\mathrm{d}x$, $B=\displaystyle\int_{0}^{1}g(x)\,\mathrm{d}x$. Ta có hệ $A+B=8$, $2A+B=10$, suy ra $A=2$, $B=6$. Vậy giá trị cần tìm là $-4$.}]{Cho các hàm số $f(x)$, $g(x)$ liên tục trên $\mathbb{R}$ thoả mãn $\displaystyle\int_{0}^{1}\left[f(x)+g(x)\right]\,\mathrm{d}x=8$ và $\displaystyle\int_{0}^{1}\left[2f(x)+g(x)\right]\,\mathrm{d}x=10$. Giá trị của $\displaystyle\int_{0}^{1}\left[f(x)-g(x)\right]\,\mathrm{d}x$ bằng}
  {$-4$}
  {$2$}{$4$}{$-2$}

\cauTN[id=C4-B12-D01-VD-029,muc=VD,loigiai={$\displaystyle\int_{1}^{3}f(x)\,\mathrm{d}x=-6-(-5)=-1$; $\displaystyle\int_{1}^{3}-6x^2\,\mathrm{d}x=-52$. Vậy $\displaystyle\int_{1}^{3}\left[2f(x)+6x^{2}\right]\,\mathrm{d}x=2\cdot(-1)-(-52)=50$.}]{Cho hàm số $f(x)$ liên tục trên $\mathbb{R}$, $\displaystyle\int_{0}^{1}f(x)\,\mathrm{d}x=-5$ và $\displaystyle\int_{0}^{3}f(x)\,\mathrm{d}x=-6$. Giá trị của $\displaystyle\int_{1}^{3}\left[2f(x)+6x^{2}\right]\,\mathrm{d}x$ bằng}
  {$50$}
  {$30$}{$-54$}{$51$}

\cauTN[id=C4-B12-D01-VD-030,muc=VD,loigiai={Đặt $A=\displaystyle\int_{0}^{2}f(x)\,\mathrm{d}x$, $B=\displaystyle\int_{0}^{2}g(x)\,\mathrm{d}x$. Ta có hệ $A+B=2$, $3A+B=10$, suy ra $A=4$, $B=-2$. Vậy giá trị cần tìm là $6$.}]{Cho các hàm số $f(x)$, $g(x)$ liên tục trên $\mathbb{R}$ thoả mãn $\displaystyle\int_{0}^{2}\left[f(x)+g(x)\right]\,\mathrm{d}x=2$ và $\displaystyle\int_{0}^{2}\left[3f(x)+g(x)\right]\,\mathrm{d}x=10$. Giá trị của $\displaystyle\int_{0}^{2}\left[f(x)-g(x)\right]\,\mathrm{d}x$ bằng}
  {$6$}
  {$4$}{$-6$}{$-8$}

\cauTN[id=C4-B12-D01-VD-031,muc=VD,loigiai={Đặt $A=\displaystyle\int_{0}^{3}f(x)\,\mathrm{d}x$, $B=\displaystyle\int_{0}^{3}g(x)\,\mathrm{d}x$. Ta có hệ $A+B=4$, $2A-B=14$, suy ra $A=6$, $B=-2$. Vậy giá trị cần tìm là $6$.}]{Cho các hàm số $f(x)$, $g(x)$ liên tục trên $\mathbb{R}$ thoả mãn $\displaystyle\int_{0}^{3}\left[f(x)+g(x)\right]\,\mathrm{d}x=4$ và $\displaystyle\int_{0}^{3}\left[2f(x)-g(x)\right]\,\mathrm{d}x=14$. Giá trị của $\displaystyle\int_{0}^{3}f(x)\,\mathrm{d}x$ bằng}
  {$6$}
  {$-2$}{$-6$}{$-10$}

\cauTN[id=C4-B12-D01-VD-032,muc=VD,loigiai={$\displaystyle\int_{1}^{4}\left[2f'(x)+3x\right]\,\mathrm{d}x=2\left[f(4)-f(1)\right]+\displaystyle\int_{1}^{4}3x\,\mathrm{d}x=2\cdot4+\dfrac{45}{2}=\dfrac{61}{2}$.}]{Cho hàm số $f(x)$ có đạo hàm liên tục trên $\mathbb{R}$, $f(1)=3$ và $f(4)=7$. Giá trị của $\displaystyle\int_{1}^{4}\left[2f'(x)+3x\right]\,\mathrm{d}x$ bằng}
  {$\dfrac{61}{2}$}
  {$\dfrac{53}{2}$}{$17$}{$\dfrac{73}{2}$}

\cauTN[id=C4-B12-D01-VD-033,muc=VD,loigiai={Đặt $A=\displaystyle\int_{0}^{4}f(x)\,\mathrm{d}x$, $B=\displaystyle\int_{0}^{4}g(x)\,\mathrm{d}x$. Ta có hệ $A+B=-5$, $3A+B=-9$, suy ra $A=-2$, $B=-3$. Vậy giá trị cần tìm là $-3$.}]{Cho các hàm số $f(x)$, $g(x)$ liên tục trên $\mathbb{R}$ thoả mãn $\displaystyle\int_{0}^{4}\left[f(x)+g(x)\right]\,\mathrm{d}x=-5$ và $\displaystyle\int_{0}^{4}\left[3f(x)+g(x)\right]\,\mathrm{d}x=-9$. Giá trị của $\displaystyle\int_{0}^{4}g(x)\,\mathrm{d}x$ bằng}
  {$-3$}
  {$-2$}{$3$}{$4$}

\cauTN[id=C4-B12-D01-VD-034,muc=VD,loigiai={Đặt $A=\displaystyle\int_{0}^{4}f(x)\,\mathrm{d}x$, $B=\displaystyle\int_{0}^{4}g(x)\,\mathrm{d}x$. Ta có hệ $A+B=-2$, $2A+B=1$, suy ra $A=3$, $B=-5$. Vậy giá trị cần tìm là $3$.}]{Cho các hàm số $f(x)$, $g(x)$ liên tục trên $\mathbb{R}$ thoả mãn $\displaystyle\int_{0}^{4}\left[f(x)+g(x)\right]\,\mathrm{d}x=-2$ và $\displaystyle\int_{0}^{4}\left[2f(x)+g(x)\right]\,\mathrm{d}x=1$. Giá trị của $\displaystyle\int_{0}^{4}f(x)\,\mathrm{d}x$ bằng}
  {$3$}
  {$-5$}{$-3$}{$-2$}

\cauTN[id=C4-B12-D01-VD-035,muc=VD,loigiai={$\displaystyle\int_{3}^{5}f(x)\,\mathrm{d}x=3-(8)=-5$; $\displaystyle\int_{3}^{5}6x^2\,\mathrm{d}x=196$. Vậy $\displaystyle\int_{3}^{5}\left[2f(x)-6x^{2}\right]\,\mathrm{d}x=2\cdot(-5)-(196)=-206$.}]{Cho hàm số $f(x)$ liên tục trên $\mathbb{R}$, $\displaystyle\int_{0}^{3}f(x)\,\mathrm{d}x=8$ và $\displaystyle\int_{0}^{5}f(x)\,\mathrm{d}x=3$. Giá trị của $\displaystyle\int_{3}^{5}\left[2f(x)-6x^{2}\right]\,\mathrm{d}x$ bằng}
  {$-206$}
  {$-174$}{$186$}{$-201$}

\cauTN[id=C4-B12-D01-VD-036,muc=VD,loigiai={Đặt $A=\displaystyle\int_{0}^{2}f(x)\,\mathrm{d}x$, $B=\displaystyle\int_{0}^{2}g(x)\,\mathrm{d}x$. Ta có hệ $A+B=-3$, $3A+B=-7$, suy ra $A=-2$, $B=-1$. Vậy giá trị cần tìm là $-2$.}]{Cho các hàm số $f(x)$, $g(x)$ liên tục trên $\mathbb{R}$ thoả mãn $\displaystyle\int_{0}^{2}\left[f(x)+g(x)\right]\,\mathrm{d}x=-3$ và $\displaystyle\int_{0}^{2}\left[3f(x)+g(x)\right]\,\mathrm{d}x=-7$. Giá trị của $\displaystyle\int_{0}^{2}f(x)\,\mathrm{d}x$ bằng}
  {$-2$}
  {$-1$}{$2$}{$4$}

\cauDS[id=C4-B12-D01-VD-037,muc=VD,loigiai={Đặt $A=\displaystyle\int_{0}^{1}f(x)\,\mathrm{d}x$, $B=\displaystyle\int_{0}^{1}g(x)\,\mathrm{d}x$: $A+B=7$, $2A-2B=2$ nên $A=4$, $B=3$. a) Đúng. b) Sai: đúng là $\dfrac{9}{2}$. c) Đúng. d) Sai: đúng là $1$.}]{Cho các hàm số $f(x)$, $g(x)$ liên tục trên $\mathbb{R}$ thoả mãn $\displaystyle\int_{0}^{1}\left[f(x)+g(x)\right]\,\mathrm{d}x=7$ và $\displaystyle\int_{0}^{1}\left[2f(x)-2g(x)\right]\,\mathrm{d}x=2$. Xét tính đúng sai của các mệnh đề sau:}{%
  \y{$\displaystyle\int_{0}^{1}g(x)\,\mathrm{d}x=3$.}{d}
  \y{$\displaystyle\int_{0}^{1}\left[f(x)+x\right]\,\mathrm{d}x=5$.}{s}
  \y{$\displaystyle\int_{0}^{1}\left[f(x)+x\right]\,\mathrm{d}x=\dfrac{9}{2}$.}{d}
  \y{$\displaystyle\int_{0}^{1}\left[f(x)-g(x)\right]\,\mathrm{d}x=-1$.}{s}}

\cauDS[id=C4-B12-D01-VD-038,muc=VD,loigiai={Đặt $A=\displaystyle\int_{0}^{3}f(x)\,\mathrm{d}x$, $B=\displaystyle\int_{0}^{3}g(x)\,\mathrm{d}x$: $A+B=-5$, $3A-2B=0$ nên $A=-2$, $B=-3$. a) Đúng. b) Đúng. c) Sai: đúng là $1$. d) Đúng.}]{Cho các hàm số $f(x)$, $g(x)$ liên tục trên $\mathbb{R}$ thoả mãn $\displaystyle\int_{0}^{3}\left[f(x)+g(x)\right]\,\mathrm{d}x=-5$ và $\displaystyle\int_{0}^{3}\left[3f(x)-2g(x)\right]\,\mathrm{d}x=0$. Xét tính đúng sai của các mệnh đề sau:}{%
  \y{$\displaystyle\int_{0}^{3}\left[f(x)+x\right]\,\mathrm{d}x=\dfrac{5}{2}$.}{d}
  \y{$\displaystyle\int_{0}^{3}\left[f(x)-g(x)\right]\,\mathrm{d}x=1$.}{d}
  \y{$\displaystyle\int_{0}^{3}\left[f(x)-g(x)\right]\,\mathrm{d}x=-1$.}{s}
  \y{$\displaystyle\int_{3}^{0}\left[f(x)+g(x)\right]\,\mathrm{d}x=5$.}{d}}

\cauDS[id=C4-B12-D01-VD-039,muc=VD,loigiai={Đặt $A=\displaystyle\int_{0}^{4}f(x)\,\mathrm{d}x$, $B=\displaystyle\int_{0}^{4}g(x)\,\mathrm{d}x$: $A+B=-7$, $3A-B=-9$ nên $A=-4$, $B=-3$. a) Sai: đúng là $-1$. b) Sai: đúng là $4$. c) Đúng. d) Sai: đúng là $-4$.}]{Cho các hàm số $f(x)$, $g(x)$ liên tục trên $\mathbb{R}$ thoả mãn $\displaystyle\int_{0}^{4}\left[f(x)+g(x)\right]\,\mathrm{d}x=-7$ và $\displaystyle\int_{0}^{4}\left[3f(x)-g(x)\right]\,\mathrm{d}x=-9$. Xét tính đúng sai của các mệnh đề sau:}{%
  \y{$\displaystyle\int_{0}^{4}\left[f(x)-g(x)\right]\,\mathrm{d}x=1$.}{s}
  \y{$\displaystyle\int_{0}^{4}\left[f(x)+x\right]\,\mathrm{d}x=12$.}{s}
  \y{$\displaystyle\int_{0}^{4}\left[f(x)+x\right]\,\mathrm{d}x=4$.}{d}
  \y{$\displaystyle\int_{0}^{4}f(x)\,\mathrm{d}x=-3$.}{s}}

\cauDS[id=C4-B12-D01-VD-040,muc=VD,loigiai={Đặt $A=\displaystyle\int_{0}^{1}f(x)\,\mathrm{d}x$, $B=\displaystyle\int_{0}^{1}g(x)\,\mathrm{d}x$: $A+B=-10$, $2A-2B=0$ nên $A=-5$, $B=-5$. a) Đúng. b) Sai: đúng là $-\dfrac{9}{2}$. c) Đúng. d) Đúng.}]{Cho các hàm số $f(x)$, $g(x)$ liên tục trên $\mathbb{R}$ thoả mãn $\displaystyle\int_{0}^{1}\left[f(x)+g(x)\right]\,\mathrm{d}x=-10$ và $\displaystyle\int_{0}^{1}\left[2f(x)-2g(x)\right]\,\mathrm{d}x=0$. Xét tính đúng sai của các mệnh đề sau:}{%
  \y{$\displaystyle\int_{0}^{1}g(x)\,\mathrm{d}x=-5$.}{d}
  \y{$\displaystyle\int_{0}^{1}\left[f(x)+x\right]\,\mathrm{d}x=-4$.}{s}
  \y{$\displaystyle\int_{0}^{1}f(x)\,\mathrm{d}x=-5$.}{d}
  \y{$\displaystyle\int_{0}^{1}\left[f(x)+x\right]\,\mathrm{d}x=-\dfrac{9}{2}$.}{d}}

\cauDS[id=C4-B12-D01-VD-041,muc=VD,loigiai={Đặt $A=\displaystyle\int_{0}^{4}f(x)\,\mathrm{d}x$, $B=\displaystyle\int_{0}^{4}g(x)\,\mathrm{d}x$: $A+B=10$, $3A-B=14$ nên $A=6$, $B=4$. a) Sai: đúng là $2$. b) Đúng. c) Đúng. d) Đúng.}]{Cho các hàm số $f(x)$, $g(x)$ liên tục trên $\mathbb{R}$ thoả mãn $\displaystyle\int_{0}^{4}\left[f(x)+g(x)\right]\,\mathrm{d}x=10$ và $\displaystyle\int_{0}^{4}\left[3f(x)-g(x)\right]\,\mathrm{d}x=14$. Xét tính đúng sai của các mệnh đề sau:}{%
  \y{$\displaystyle\int_{0}^{4}\left[f(x)-g(x)\right]\,\mathrm{d}x=-2$.}{s}
  \y{$\displaystyle\int_{0}^{4}f(x)\,\mathrm{d}x=6$.}{d}
  \y{$\displaystyle\int_{0}^{4}\left[f(x)-g(x)\right]\,\mathrm{d}x=2$.}{d}
  \y{$\displaystyle\int_{0}^{4}g(x)\,\mathrm{d}x=4$.}{d}}

\cauDS[id=C4-B12-D01-VD-042,muc=VD,loigiai={Đặt $A=\displaystyle\int_{0}^{4}f(x)\,\mathrm{d}x$, $B=\displaystyle\int_{0}^{4}g(x)\,\mathrm{d}x$: $A+B=12$, $3A-B=12$ nên $A=6$, $B=6$. a) Đúng: đúng là $6$. b) Đúng. c) Sai: đúng là $14$. d) Đúng.}]{Cho các hàm số $f(x)$, $g(x)$ liên tục trên $\mathbb{R}$ thoả mãn $\displaystyle\int_{0}^{4}\left[f(x)+g(x)\right]\,\mathrm{d}x=12$ và $\displaystyle\int_{0}^{4}\left[3f(x)-g(x)\right]\,\mathrm{d}x=12$. Xét tính đúng sai của các mệnh đề sau:}{%
  \y{$\displaystyle\int_{0}^{4}f(x)\,\mathrm{d}x=6$.}{d}
  \y{$\displaystyle\int_{0}^{4}g(x)\,\mathrm{d}x=6$.}{d}
  \y{$\displaystyle\int_{0}^{4}\left[f(x)+x\right]\,\mathrm{d}x=22$.}{s}
  \y{$\displaystyle\int_{0}^{4}\left[f(x)+x\right]\,\mathrm{d}x=14$.}{d}}

\cauDS[id=C4-B12-D01-VD-043,muc=VD,loigiai={Đặt $A=\displaystyle\int_{0}^{2}f(x)\,\mathrm{d}x$, $B=\displaystyle\int_{0}^{2}g(x)\,\mathrm{d}x$: $A+B=-8$, $2A-B=-10$ nên $A=-6$, $B=-2$. a) Sai: đúng là $-4$. b) Đúng. c) Sai: đúng là $-6$. d) Sai: đúng là $-4$.}]{Cho các hàm số $f(x)$, $g(x)$ liên tục trên $\mathbb{R}$ thoả mãn $\displaystyle\int_{0}^{2}\left[f(x)+g(x)\right]\,\mathrm{d}x=-8$ và $\displaystyle\int_{0}^{2}\left[2f(x)-g(x)\right]\,\mathrm{d}x=-10$. Xét tính đúng sai của các mệnh đề sau:}{%
  \y{$\displaystyle\int_{0}^{2}\left[f(x)+x\right]\,\mathrm{d}x=-2$.}{s}
  \y{$\displaystyle\int_{0}^{2}\left[f(x)+x\right]\,\mathrm{d}x=-4$.}{d}
  \y{$\displaystyle\int_{0}^{2}f(x)\,\mathrm{d}x=-2$.}{s}
  \y{$\displaystyle\int_{0}^{2}\left[f(x)-g(x)\right]\,\mathrm{d}x=4$.}{s}}

\cauDS[id=C4-B12-D01-VD-044,muc=VD,loigiai={Đặt $A=\displaystyle\int_{0}^{1}f(x)\,\mathrm{d}x$, $B=\displaystyle\int_{0}^{1}g(x)\,\mathrm{d}x$: $A+B=-9$, $2A-2B=-6$ nên $A=-6$, $B=-3$. a) Đúng. b) Sai: đúng là $-6$. c) Đúng. d) Đúng.}]{Cho các hàm số $f(x)$, $g(x)$ liên tục trên $\mathbb{R}$ thoả mãn $\displaystyle\int_{0}^{1}\left[f(x)+g(x)\right]\,\mathrm{d}x=-9$ và $\displaystyle\int_{0}^{1}\left[2f(x)-2g(x)\right]\,\mathrm{d}x=-6$. Xét tính đúng sai của các mệnh đề sau:}{%
  \y{$\displaystyle\int_{1}^{0}\left[f(x)+g(x)\right]\,\mathrm{d}x=9$.}{d}
  \y{$\displaystyle\int_{0}^{1}f(x)\,\mathrm{d}x=-3$.}{s}
  \y{$\displaystyle\int_{0}^{1}\left[f(x)-g(x)\right]\,\mathrm{d}x=-3$.}{d}
  \y{$\displaystyle\int_{0}^{1}\left[f(x)+x\right]\,\mathrm{d}x=-\dfrac{11}{2}$.}{d}}

\cauDS[id=C4-B12-D01-VD-045,muc=VD,loigiai={Đặt $A=\displaystyle\int_{0}^{1}f(x)\,\mathrm{d}x$, $B=\displaystyle\int_{0}^{1}g(x)\,\mathrm{d}x$: $A+B=-1$, $3A-B=-7$ nên $A=-2$, $B=1$. a) Đúng. b) Sai: đúng là $-\dfrac{3}{2}$. c) Sai: đúng là $-3$. d) Đúng.}]{Cho các hàm số $f(x)$, $g(x)$ liên tục trên $\mathbb{R}$ thoả mãn $\displaystyle\int_{0}^{1}\left[f(x)+g(x)\right]\,\mathrm{d}x=-1$ và $\displaystyle\int_{0}^{1}\left[3f(x)-g(x)\right]\,\mathrm{d}x=-7$. Xét tính đúng sai của các mệnh đề sau:}{%
  \y{$\displaystyle\int_{0}^{1}\left[f(x)-g(x)\right]\,\mathrm{d}x=-3$.}{d}
  \y{$\displaystyle\int_{0}^{1}\left[f(x)+x\right]\,\mathrm{d}x=-1$.}{s}
  \y{$\displaystyle\int_{0}^{1}\left[f(x)-g(x)\right]\,\mathrm{d}x=3$.}{s}
  \y{$\displaystyle\int_{0}^{1}f(x)\,\mathrm{d}x=-2$.}{d}}

\cauDS[id=C4-B12-D01-VD-046,muc=VD,loigiai={Đặt $A=\displaystyle\int_{0}^{4}f(x)\,\mathrm{d}x$, $B=\displaystyle\int_{0}^{4}g(x)\,\mathrm{d}x$: $A+B=2$, $3A-B=-18$ nên $A=-4$, $B=6$. a) Sai: đúng là $-10$. b) Sai: đúng là $-4$. c) Sai: đúng là $4$. d) Đúng.}]{Cho các hàm số $f(x)$, $g(x)$ liên tục trên $\mathbb{R}$ thoả mãn $\displaystyle\int_{0}^{4}\left[f(x)+g(x)\right]\,\mathrm{d}x=2$ và $\displaystyle\int_{0}^{4}\left[3f(x)-g(x)\right]\,\mathrm{d}x=-18$. Xét tính đúng sai của các mệnh đề sau:}{%
  \y{$\displaystyle\int_{0}^{4}\left[f(x)-g(x)\right]\,\mathrm{d}x=10$.}{s}
  \y{$\displaystyle\int_{0}^{4}f(x)\,\mathrm{d}x=6$.}{s}
  \y{$\displaystyle\int_{0}^{4}\left[f(x)+x\right]\,\mathrm{d}x=12$.}{s}
  \y{$\displaystyle\int_{4}^{0}\left[f(x)+g(x)\right]\,\mathrm{d}x=-2$.}{d}}

\cauDS[id=C4-B12-D01-VD-047,muc=VD,loigiai={Đặt $A=\displaystyle\int_{0}^{3}f(x)\,\mathrm{d}x$, $B=\displaystyle\int_{0}^{3}g(x)\,\mathrm{d}x$: $A+B=-7$, $3A-2B=-6$ nên $A=-4$, $B=-3$. a) Đúng. b) Đúng. c) Sai: đúng là $\dfrac{1}{2}$. d) Sai: đúng là $-4$.}]{Cho các hàm số $f(x)$, $g(x)$ liên tục trên $\mathbb{R}$ thoả mãn $\displaystyle\int_{0}^{3}\left[f(x)+g(x)\right]\,\mathrm{d}x=-7$ và $\displaystyle\int_{0}^{3}\left[3f(x)-2g(x)\right]\,\mathrm{d}x=-6$. Xét tính đúng sai của các mệnh đề sau:}{%
  \y{$\displaystyle\int_{0}^{3}\left[f(x)-g(x)\right]\,\mathrm{d}x=-1$.}{d}
  \y{$\displaystyle\int_{0}^{3}\left[f(x)+x\right]\,\mathrm{d}x=\dfrac{1}{2}$.}{d}
  \y{$\displaystyle\int_{0}^{3}\left[f(x)+x\right]\,\mathrm{d}x=5$.}{s}
  \y{$\displaystyle\int_{0}^{3}f(x)\,\mathrm{d}x=-3$.}{s}}

\cauDS[id=C4-B12-D01-VD-048,muc=VD,loigiai={Đặt $A=\displaystyle\int_{0}^{4}f(x)\,\mathrm{d}x$, $B=\displaystyle\int_{0}^{4}g(x)\,\mathrm{d}x$: $A+B=-2$, $2A-B=14$ nên $A=4$, $B=-6$. a) Đúng. b) Sai: đúng là $10$. c) Đúng. d) Sai: đúng là $12$.}]{Cho các hàm số $f(x)$, $g(x)$ liên tục trên $\mathbb{R}$ thoả mãn $\displaystyle\int_{0}^{4}\left[f(x)+g(x)\right]\,\mathrm{d}x=-2$ và $\displaystyle\int_{0}^{4}\left[2f(x)-g(x)\right]\,\mathrm{d}x=14$. Xét tính đúng sai của các mệnh đề sau:}{%
  \y{$\displaystyle\int_{4}^{0}\left[f(x)+g(x)\right]\,\mathrm{d}x=2$.}{d}
  \y{$\displaystyle\int_{0}^{4}\left[f(x)-g(x)\right]\,\mathrm{d}x=-10$.}{s}
  \y{$\displaystyle\int_{0}^{4}f(x)\,\mathrm{d}x=4$.}{d}
  \y{$\displaystyle\int_{0}^{4}\left[f(x)+x\right]\,\mathrm{d}x=20$.}{s}}

\cauTLN[id=C4-B12-D01-VD-049,muc=VD,loigiai={$\displaystyle\int_{1}^{3}\left[2f'(x)+2x\right]\,\mathrm{d}x=2\left[f(3)-f(1)\right]+\displaystyle\int_{1}^{3}2x\,\mathrm{d}x=2\cdot-2+8=4$.}]{Cho hàm số $f(x)$ có đạo hàm liên tục trên $\mathbb{R}$, $f(1)=1$ và $f(3)=-1$. Tính $\displaystyle\int_{1}^{3}\left[2f'(x)+2x\right]\,\mathrm{d}x$.}{4}

\cauTLN[id=C4-B12-D01-VD-050,muc=VD,loigiai={Đặt $A=\displaystyle\int_{0}^{1}f(x)\,\mathrm{d}x$, $B=\displaystyle\int_{0}^{1}g(x)\,\mathrm{d}x$. Ta có hệ $A+B=-1$, $2A-2B=-18$, suy ra $A=-5$, $B=4$. Vậy giá trị cần tìm là $-5$.}]{Cho các hàm số $f(x)$, $g(x)$ liên tục trên $\mathbb{R}$ thoả mãn $\displaystyle\int_{0}^{1}\left[f(x)+g(x)\right]\,\mathrm{d}x=-1$ và $\displaystyle\int_{0}^{1}\left[2f(x)-2g(x)\right]\,\mathrm{d}x=-18$. Tính $\displaystyle\int_{0}^{1}f(x)\,\mathrm{d}x$.}{-5}

\cauTLN[id=C4-B12-D01-VD-051,muc=VD,loigiai={Đặt $A=\displaystyle\int_{0}^{4}f(x)\,\mathrm{d}x$, $B=\displaystyle\int_{0}^{4}g(x)\,\mathrm{d}x$. Ta có hệ $A+B=-3$, $2A+B=-1$, suy ra $A=2$, $B=-5$. Vậy giá trị cần tìm là $2$.}]{Cho các hàm số $f(x)$, $g(x)$ liên tục trên $\mathbb{R}$ thoả mãn $\displaystyle\int_{0}^{4}\left[f(x)+g(x)\right]\,\mathrm{d}x=-3$ và $\displaystyle\int_{0}^{4}\left[2f(x)+g(x)\right]\,\mathrm{d}x=-1$. Tính $\displaystyle\int_{0}^{4}f(x)\,\mathrm{d}x$.}{2}

\cauTLN[id=C4-B12-D01-VD-052,muc=VD,loigiai={$\displaystyle\int_{0}^{3}\left[2f'(x)-x\right]\,\mathrm{d}x=2\left[f(3)-f(0)\right]+\displaystyle\int_{0}^{3}-x\,\mathrm{d}x=2\cdot-2-\dfrac{9}{2}=-\dfrac{17}{2}$.}]{Cho hàm số $f(x)$ có đạo hàm liên tục trên $\mathbb{R}$, $f(0)=1$ và $f(3)=-1$. Tính $\displaystyle\int_{0}^{3}\left[2f'(x)-x\right]\,\mathrm{d}x$.}{-8,5}

\cauTLN[id=C4-B12-D01-VD-053,muc=VD,loigiai={$\displaystyle\int_{1}^{3}f(x)\,\mathrm{d}x=7-(2)=5$; $\displaystyle\int_{1}^{3}6x^2\,\mathrm{d}x=52$. Vậy $\displaystyle\int_{1}^{3}\left[2f(x)-6x^{2}\right]\,\mathrm{d}x=2\cdot(5)-(52)=-42$.}]{Cho hàm số $f(x)$ liên tục trên $\mathbb{R}$, $\displaystyle\int_{0}^{1}f(x)\,\mathrm{d}x=2$ và $\displaystyle\int_{0}^{3}f(x)\,\mathrm{d}x=7$. Tính $\displaystyle\int_{1}^{3}\left[2f(x)-6x^{2}\right]\,\mathrm{d}x$.}{-42}

\cauTLN[id=C4-B12-D01-VD-054,muc=VD,loigiai={$\displaystyle\int_{-1}^{1}\left[2f'(x)-x\right]\,\mathrm{d}x=2\left[f(1)-f(-1)\right]+\displaystyle\int_{-1}^{1}-x\,\mathrm{d}x=2\cdot-6+0=-12$.}]{Cho hàm số $f(x)$ có đạo hàm liên tục trên $\mathbb{R}$, $f(-1)=5$ và $f(1)=-1$. Tính $\displaystyle\int_{-1}^{1}\left[2f'(x)-x\right]\,\mathrm{d}x$.}{-12}

\cauTLN[id=C4-B12-D01-VD-055,muc=VD,loigiai={$\displaystyle\int_{1}^{4}f(x)\,\mathrm{d}x=-6-(-5)=-1$; $\displaystyle\int_{1}^{4}-6x^2\,\mathrm{d}x=-126$. Vậy $\displaystyle\int_{1}^{4}\left[2f(x)+6x^{2}\right]\,\mathrm{d}x=2\cdot(-1)-(-126)=124$.}]{Cho hàm số $f(x)$ liên tục trên $\mathbb{R}$, $\displaystyle\int_{0}^{1}f(x)\,\mathrm{d}x=-5$ và $\displaystyle\int_{0}^{4}f(x)\,\mathrm{d}x=-6$. Tính $\displaystyle\int_{1}^{4}\left[2f(x)+6x^{2}\right]\,\mathrm{d}x$.}{124}

\cauTLN[id=C4-B12-D01-VD-056,muc=VD,loigiai={$\displaystyle\int_{-1}^{1}\left[2f'(x)+2x\right]\,\mathrm{d}x=2\left[f(1)-f(-1)\right]+\displaystyle\int_{-1}^{1}2x\,\mathrm{d}x=2\cdot-2+0=-4$.}]{Cho hàm số $f(x)$ có đạo hàm liên tục trên $\mathbb{R}$, $f(-1)=5$ và $f(1)=3$. Tính $\displaystyle\int_{-1}^{1}\left[2f'(x)+2x\right]\,\mathrm{d}x$.}{-4}

\cauTLN[id=C4-B12-D01-VD-057,muc=VD,loigiai={$\displaystyle\int_{1}^{2}f(x)\,\mathrm{d}x=-6-(-5)=-1$; $\displaystyle\int_{1}^{2}-3x^2\,\mathrm{d}x=-7$. Vậy $\displaystyle\int_{1}^{2}\left[2f(x)+3x^{2}\right]\,\mathrm{d}x=2\cdot(-1)-(-7)=5$.}]{Cho hàm số $f(x)$ liên tục trên $\mathbb{R}$, $\displaystyle\int_{0}^{1}f(x)\,\mathrm{d}x=-5$ và $\displaystyle\int_{0}^{2}f(x)\,\mathrm{d}x=-6$. Tính $\displaystyle\int_{1}^{2}\left[2f(x)+3x^{2}\right]\,\mathrm{d}x$.}{5}

\cauTLN[id=C4-B12-D01-VD-058,muc=VD,loigiai={$\displaystyle\int_{1}^{3}f(x)\,\mathrm{d}x=4-(6)=-2$; $\displaystyle\int_{1}^{3}-6x^2\,\mathrm{d}x=-52$. Vậy $\displaystyle\int_{1}^{3}\left[2f(x)+6x^{2}\right]\,\mathrm{d}x=2\cdot(-2)-(-52)=48$.}]{Cho hàm số $f(x)$ liên tục trên $\mathbb{R}$, $\displaystyle\int_{0}^{1}f(x)\,\mathrm{d}x=6$ và $\displaystyle\int_{0}^{3}f(x)\,\mathrm{d}x=4$. Tính $\displaystyle\int_{1}^{3}\left[2f(x)+6x^{2}\right]\,\mathrm{d}x$.}{48}

\cauTLN[id=C4-B12-D01-VD-059,muc=VD,loigiai={$\displaystyle\int_{2}^{3}f(x)\,\mathrm{d}x=-3-(7)=-10$; $\displaystyle\int_{2}^{3}6x^2\,\mathrm{d}x=38$. Vậy $\displaystyle\int_{2}^{3}\left[2f(x)-6x^{2}\right]\,\mathrm{d}x=2\cdot(-10)-(38)=-58$.}]{Cho hàm số $f(x)$ liên tục trên $\mathbb{R}$, $\displaystyle\int_{0}^{2}f(x)\,\mathrm{d}x=7$ và $\displaystyle\int_{0}^{3}f(x)\,\mathrm{d}x=-3$. Tính $\displaystyle\int_{2}^{3}\left[2f(x)-6x^{2}\right]\,\mathrm{d}x$.}{-58}

\cauTLN[id=C4-B12-D01-VD-060,muc=VD,loigiai={$\displaystyle\int_{0}^{3}\left[2f'(x)-x\right]\,\mathrm{d}x=2\left[f(3)-f(0)\right]+\displaystyle\int_{0}^{3}-x\,\mathrm{d}x=2\cdot-1-\dfrac{9}{2}=-\dfrac{13}{2}$.}]{Cho hàm số $f(x)$ có đạo hàm liên tục trên $\mathbb{R}$, $f(0)=2$ và $f(3)=1$. Tính $\displaystyle\int_{0}^{3}\left[2f'(x)-x\right]\,\mathrm{d}x$.}{-6,5}
